\documentclass[11pt]{article}
\usepackage{iftex}
\usepackage[T1]{fontenc}
\ifPDFTeX\usepackage[utf8]{inputenc}\fi
\usepackage{lmodern,amsmath,amssymb,amsthm,geometry,array,longtable,booktabs,hyperref}
\geometry{margin=25mm}
\hypersetup{colorlinks=true,linkcolor=blue,citecolor=blue,urlcolor=blue,
pdftitle={Extremal relations and Betti numbers of four-generated numerical semigroups},
pdfauthor={Keisuke Sato},
pdfsubject={Numerical semigroups and total Betti numbers},
pdfkeywords={numerical semigroups, minimal presentations, Betti numbers, multiplicity, Apery sets}}
\allowdisplaybreaks
\setlength{\emergencystretch}{3em}
\newtheorem{theorem}{Theorem}[section]
\newtheorem{lemma}[theorem]{Lemma}
\newtheorem{proposition}[theorem]{Proposition}
\newtheorem{corollary}[theorem]{Corollary}
\theoremstyle{definition}
\newtheorem{example}[theorem]{Example}
\theoremstyle{remark}
\newtheorem{remark}[theorem]{Remark}
\theoremstyle{plain}
\DeclareMathOperator{\Ap}{Ap}
\DeclareMathOperator{\type}{type}
\DeclareMathOperator{\Soc}{Soc}
\DeclareMathOperator{\edim}{edim}
\ifdefined\DeclareUnicodeCharacter
\DeclareUnicodeCharacter{0131}{\i}
\DeclareUnicodeCharacter{011F}{\u{g}}
\DeclareUnicodeCharacter{0161}{\v{s}}
\fi
\title{Extremal relations and Betti numbers of\\four-generated numerical semigroups}
\author{Keisuke Sato\\Independent researcher, Japan\\
\href{mailto:stksk310@gmail.com}{stksk310@gmail.com}\\
ORCID: \href{https://orcid.org/0009-0005-4230-0268}{0009-0005-4230-0268}}
\date{19 November 2026}
\begin{document}
\maketitle
\begin{abstract}
For numerical semigroups minimally generated by four integers, we prove that
the largest cardinality of a minimal presentation at multiplicity $m\ge4$ is
$T(m)=2+\lfloor\sqrt{4m}\rfloor$.
The new ingredient is a universal upper bound: canonical Ap\'ery
representatives form a finite downset in $\mathbb N^3$, and its minimal
excluded exponent vectors number at most $T(m)$. Published constructions then give attainment
and the complete attainable spectrum, proving Conjecture~7.1 of
Elmac\i{}o\u{g}lu, Hilmer, O'Neill, Okandan and Park-Kaufmann.
We also restrict each of the maximal type and second total Betti number
to at most three consecutive integer possibilities, and characterize
equality with the Macaulay bound. The main results are formally verified
in Lean~4.
\end{abstract}
\noindent\textbf{Keywords:} numerical semigroup; minimal presentation;
Betti numbers; multiplicity; Ap\'ery set.\par
\smallskip
\noindent\textbf{MSC 2020:} 20M14, 13D02, 13F65.

\section{Introduction}
Let $S=\langle m,n_1,n_2,n_3\rangle$ be a numerical semigroup whose
displayed generators are minimal, with $m=\min(S\setminus\{0\})$.
Write $\eta(S)$ for the common cardinality of its minimal presentations,
equivalently the minimal number of generators of its defining toric ideal.
Herzog showed that $\eta(S)\le3$ for three-generated numerical
semigroups \cite{Herzog70} (see also~\cite{HerzogStamate14}). Bresinsky showed that
these cardinalities are unbounded already for
four-generated numerical semigroups as the generators vary
\cite{Bresinsky75}. Fixing the multiplicity therefore gives a natural
extremal problem. Recent bounds and questions are developed by
Moscariello and Sammartano \cite{MS25,MSSurvey25}.

\begin{theorem}\label{thm:main}
For every integer $m\ge4$,
\[
\max\{\eta(S):\edim S=4,\ m(S)=m\}
=T(m):=2+\lfloor\sqrt{4m}\rfloor.
\]
\end{theorem}

The contribution here is the upper bound for every four-generated
numerical semigroup. Elmac\i{}o\u{g}lu, Hilmer, O'Neill, Okandan and
Park-Kaufmann \cite[Theorem~6.3]{EHOOPK24} already supply attainment:
each pair $q\ge6$, $4m\ge(q-2)^2$ occurs. Together with their results
for $q=3,4,5$, the upper bound determines the complete spectrum
(Corollary~\ref{cor:spectrum}) and proves their Conjecture~7.1.

Moscariello and Sammartano obtain their Macaulay bound through Hilbert
functions and free resolutions \cite[Sections~2.2--2.3]{MS25}.
Our reduction also leads to a finite monomial staircase, but retains
the additional constraint that its points represent distinct cyclic
residues. This constraint is the source of the sharp estimate below.
Choose canonical Ap\'ery representatives $D\subset\mathbb N^3$, one
in each residue modulo $m$, and write $B=\partial D$ for their minimal
excluded exponent vectors. A normal-form reduction gives
$\eta(S)\le|B|$. Restrictions on the supports of boundary returns then
yield
\[
 (|B|-2)^2\le4|D|.
\]
The main work is to count disjoint axis, plane and interior resources,
including the cases with small plane-boundary counts.

Brower, McDonough and O'Neill establish support restrictions for outer
Betti elements \cite[Proposition~5.2]{BMO25} and a passage to staircase
Kunz nilsemigroups that does not decrease presentation cardinality
\cite[Theorem~6.2]{BMO25}. These results provide geometric antecedents
for the support arguments used here. Chomicz develops a geometric
description of Ap\'ery sets in embedding dimension four
\cite[Section~2.2]{Chomicz26A}; using this geometry, he also proves
$\tau(S)-11\le\eta(S)\le4\tau(S)+5$, where $\tau(S)$ is the type
\cite[Theorem~1.4]{Chomicz26B}. Those inequalities compare invariants
of an individual semigroup; our extremal problem fixes multiplicity
and maximizes presentation cardinality.

Section~\ref{sec:betti} combines the capacity bound with standard
resolution theory and the known attaining families. It gives simultaneous
Betti bounds and leaves at most three integer possibilities for each of
the maximal type and second total Betti number. The attaining families
are those of \cite{EHOOPK24}; the exact latter two extrema and all
attainable Betti tuples remain undetermined in general.

Section~\ref{sec:canonical} constructs the representatives and the
border presentation, with a small example. Section~\ref{sec:capacity}
proves the capacity inequality; its opening roadmap identifies the
case split. Section~\ref{sec:betti} gives the Betti, spectrum and
Macaulay consequences. The final section records the Lean interfaces
and archival sources.

We use $\mathbb N=\{0,1,\ldots\}$. Triple inequalities are coordinatewise,
and $[u]_+=\max(u,0)$. All other differences are integer differences.
In the shifted embedding-index notation of \cite{MS25},
$\mathcal R(3,m)$ is the maximum on the left of Theorem~\ref{thm:main},
and $\mathcal C(3,m)$ is its Macaulay upper bound, recalled in
Section~\ref{sec:macaulay}. The shifted index $3$ corresponds to
four minimal generators here.

\section{Canonical Ap\'ery representatives}
\label{sec:canonical}
Set $w(x,y,z)=n_1x+n_2y+n_3z$ and
$\sigma(q)=w(q)\bmod m$. Since $\gcd(m,n_1,n_2,n_3)=1$, the residues
of $n_1,n_2,n_3$ generate $\mathbb Z/m\mathbb Z$. They also generate it
as a monoid, so every residue is represented by a nonnegative triple.
In each residue choose the triple of least lexicographic key
\[
\operatorname{key}(q)=(w(q),x,y,z),
\]
and let $D$ be the resulting set. The order is well founded and
translation invariant. If $p+d\in D$ and $u$ has the residue of $p$,
then $\operatorname{key}(p+d)\le \operatorname{key}(u+d)$, hence $\operatorname{key}(p)\le \operatorname{key}(u)$. Thus $D$ is a
downset, and $\sigma|_D$ is a bijection onto $\mathbb Z/m\mathbb Z$.
In particular $|D|=m$.

The weights $w(D)$ form the Ap\'ery set of $S$ with respect to $m$
(see \cite{Apery46} and \cite[Chapter~2]{RG09}):
$\Ap(S,m)=\{s\in S:s-m\notin S\}$, where subtraction is in
$\mathbb Z$. Indeed they are the least semigroup elements of their
respective residues. Every $s\in S$ consequently has a unique form
$s=hm+w(q)$ with $h\in\mathbb N$ and $q\in D$.

For any nonempty finite downset define its boundary, or minimal complement, by
\[
\partial D=\{b\notin D:b-e_i\in D\text{ whenever }b_i>0\}.
\]
It is finite, since it lies in $\bigcup_i(D+e_i)$, and every outside
point dominates one of its members.

The next proposition uses the standard normal-form reduction pattern.
We include the argument to specify the chosen border relations.

\begin{proposition}\label{prop:border}
The semigroup $S$ has a presentation with at most $|\partial D|$ pairs.
Consequently $\eta(S)\le|\partial D|$.
\end{proposition}
\begin{proof}
For $b\in\partial D$, let $q_b\in D$ have the same residue.
Minimality gives
$w(b)=w(q_b)+h_bm$ with $h_b\ge0$ and $\operatorname{key}(q_b)<\operatorname{key}(b)$.
Take the relation $((0,b),(h_b,q_b))$.
If a factorization $(h,p)$ has $p\notin D$, choose $b\le p$ in the
boundary and replace it by $(h+h_b,p-b+q_b)$. This preserves value
and strictly decreases the triple key. It therefore terminates at the
unique normal form $(\ell,q)$, $q\in D$. Two factorizations of the same
element reduce to the same normal form, so these pairs generate the
factorization congruence. Removing redundant pairs gives a minimal
presentation and proves the inequality.
\end{proof}

\begin{example}\label{ex:small}
For $S=\langle4,7,9,10\rangle$, the least representatives of residues
$0,1,2,3$ are $0,9,10,7$, respectively. Their exponent vectors are
\[
D=\{0,e_1,e_2,e_3\},\qquad
\partial D=\{2e_1,2e_2,2e_3,e_1+e_2,e_1+e_3,e_2+e_3\}.
\]
Thus $|D|=4$ and $|\partial D|=6=T(4)$. The return of $2e_1$ is
$e_3$, since $14\equiv10\pmod4$, and its border relation replaces
two copies of $7$ by $4+10$. Figure~\ref{fig:small} shows all points
of $D$ and its boundary. There is one boundary point in each coordinate
plane with two positive coordinates and none with three. Thus the
profile is $(r,s,t)=(1,1,1)$ with $\epsilon=0$, and the count in
Lemma~\ref{res:A.2} gives $3+1+1+1+0=6$.
\end{example}

\begin{figure}[htbp]
\centering
\setlength{\unitlength}{1pt}
\begin{picture}(355,105)
\put(8,15){\vector(1,0){76}}\put(8,15){\vector(0,1){76}}
\put(135,15){\vector(1,0){76}}\put(135,15){\vector(0,1){76}}
\put(262,15){\vector(1,0){76}}\put(262,15){\vector(0,1){76}}
\put(87,12){$x$}\put(3,94){$y$}
\put(214,12){$x$}\put(130,94){$y$}
\put(341,12){$x$}\put(257,94){$y$}
\put(20,94){$z=0$}\put(147,94){$z=1$}\put(274,94){$z=2$}
\put(8,15){\circle*{4}}\put(38,15){\circle*{4}}\put(8,45){\circle*{4}}
\put(68,15){\circle{6}}\put(38,45){\circle{6}}\put(8,75){\circle{6}}
\put(135,15){\circle*{4}}\put(165,15){\circle{6}}\put(135,45){\circle{6}}
\put(262,15){\circle{6}}
\put(5,3){$0$}\put(35,3){$1$}\put(65,3){$2$}
\put(132,3){$0$}\put(162,3){$1$}\put(192,3){$2$}
\put(259,3){$0$}\put(289,3){$1$}\put(319,3){$2$}
\put(-5,42){$1$}\put(-5,72){$2$}
\put(122,42){$1$}\put(122,72){$2$}
\put(249,42){$1$}\put(249,72){$2$}
\end{picture}
\caption{Canonical representatives (filled dots) and minimal excluded
vectors (open circles) for $\langle4,7,9,10\rangle$, shown in horizontal slices.}
\label{fig:small}
\end{figure}

\begin{proof}[Proof of Theorem~\ref{thm:main}, assuming the capacity bound]
Theorem~\ref{thm:capacity} applied to this $D$ gives
$\eta(S)\le|\partial D|\le T(m)$.
Conversely $T(m)\ge6$ and $(T(m)-2)^2\le4m$, so
\cite[Theorem~6.3]{EHOOPK24} supplies a semigroup with
$\eta(S)=T(m)$ at every $m\ge4$.
\end{proof}

At equality $\eta(S)=T(m)$, both inequalities are equalities.
The canonical border presentation then has $T(m)$ distinct pairs
and is minimal.

\section{The cyclic capacity inequality}\label{sec:capacity}
Here $D$ is any finite downset and the additive map
$\sigma(x,y,z)=xp_1+yp_2+zp_3$ restricts to a bijection
$D\to\mathbb Z/m\mathbb Z$, with $m>0$.
Write $B=\partial D$ and $R(q)$ for the representative of $\sigma(q)$
in $D$, called its \emph{return}. The weights $p_i$ are arbitrary
elements of $\mathbb Z/m\mathbb Z$; bijectivity is required only on $D$.
A point of $D$ is called \emph{standard}. We call such a downset, together
with its residue map, a \emph{cyclic staircase}.

\begin{theorem}\label{thm:capacity}
Every such cyclic staircase satisfies
$|B|\le2+\lfloor\sqrt{4m}\rfloor$.
\end{theorem}


\paragraph{Proof strategy and notation.}
The numbers $r,s,t$ count boundary points supported on exactly the
$xy,xz,yz$ planes. In the case analysis we permute coordinates so that
$r\ge s\ge t$ and call $t$ the \emph{width}; $(r,s,t)$ is its
\emph{width profile}. A boundary point is
\emph{full-support} when all three coordinates are positive; its
uniqueness, when it exists, follows from Lemma~\ref{res:A.1}.
The first subsection develops support restrictions and disjoint
resources. The remaining subsections establish the appropriate quadratic
lower bound for $m$, according to the following table.
\begin{center}
\begin{tabular}{@{}lll@{}}\toprule
Boundary regime&Width&Closing result\\\midrule
No full-support point&all&Lemma~\ref{res:A.22}\\
Full-support point&$0$&Lemma~\ref{res:A.19}\\
Full-support point&$1$&Lemma~\ref{res:A12.1}\\
Full-support point&$2$&Theorem~\ref{res:A12.8}\\
Full-support point&$3$&Theorem~\ref{res:A3.9}\\
Full-support point&$\ge4$&Theorem~\ref{res:A.35}\\\bottomrule
\end{tabular}
\end{center}
The last row separates widths at least seven from widths four to six;
the latter leave eight explicit profiles. Their finite checks occur only
after bounds on all remaining coordinates have been proved.
The completion at the end assembles these branches and converts the
quadratic inequality to the stated square-root bound.

\begin{center}
\begin{tabular}{@{}p{.21\linewidth}p{.72\linewidth}@{}}\toprule
Notation&Meaning within this section\\\midrule
$D$, $B=\partial D$&Standard exponent vectors and minimal excluded vectors\\
$R(q)$&The unique standard vector with residue $\sigma(q)$\\
$r,s,t$&Pair-support boundary counts; ordered $r\ge s\ge t$ in the case split\\
$\delta,\delta_{st}$&The ordered gaps $r-s$ and $s-t$, respectively\\
$g=(A,B,C_0)$&The full-support boundary point, when present\\
$C$, $C^\circ$, $H$&Forced interior core, its puncturing at $g$, and the interior resource union\\
$a_i,L_i,J$&First excluded axis index, axis baseline, and joint axis gain\\
$E_{ij},F$&Extra plane points and their lower bounds\\
$Q,W,\Phi$&Rectangle, high-corner gap, and rank-cut bounds for plane excess\\
$\Pi(s,t)$&Number of points in the positive simplex used to bound $|C|$\\
$E_n,d_n$&Rectangle bound and reduced plane bound in width three\\
$\ell_{\Pi},\ell_{\mathrm{row}},\ell_{\mathrm{core}}$&Simplex lower bound, row endpoint, and core lower bound\\
$\Lambda$&A local quadratic capacity target; $T(m)$ remains the extremal function\\\bottomrule
\end{tabular}
\end{center}

\subsection{Returns and disjoint resources}
Let $a_i$ be the least positive integer with $a_ie_i\notin D$.
Thus the standard points on the $i$th axis are $0,e_i,\ldots,(a_i-1)e_i$.
\begin{lemma}[return supports]\label{res:A.1}
If $b\in B$ and $b_i>0$, then $R(b)_i=0$.
Distinct boundary points sharing a positive coordinate have different
returns. There is at most one full-support boundary point, and its return
is zero.
\end{lemma}
\begin{proof}
If $R(b)_i>0$, the standard points $b-e_i$ and $R(b)-e_i$ have the
same residue. Injectivity gives $b=R(b)$, a contradiction.
Likewise, if $b_i,c_i>0$ and $R(b)=R(c)$, injectivity at $b-e_i,c-e_i$
gives $b=c$. A full-support return must be zero, proving uniqueness.
\end{proof}

\begin{lemma}[boundary count]\label{res:A.2}
Let $r,s,t$ count the boundary points supported on exactly $xy,xz,yz$,
and let $\epsilon\in\{0,1\}$ count full-support points. Then
\[
 |B|=3+r+s+t+\epsilon.
\]
\end{lemma}
\begin{proof}
The three axis points are $a_ie_i$. Together with the three pair-support
classes and the full-support class, they partition $B$. \end{proof}

\begin{lemma}[plane ranks]\label{res:A.3}
Order the $xy$ corners as $(X_j,Y_j,0)$, $1\le j\le r$, by increasing
$X_j$. Then $Y_j$ strictly decreases,
$X_j\ge j$, $Y_j\ge r-j+1$, and $a_1,a_2\ge r+1$.
Every positive $(x,y,0)$ with $x+y\le r$ is standard.
The analogous assertions hold in the other planes.
\end{lemma}
\begin{proof}
Boundary points are incomparable. Thus the positive $X_j$ increase
strictly, the positive $Y_j$ decrease strictly, and all are below the
corresponding intercepts. This gives the rank and intercept bounds.
An outside positive point with $x+y\le r$ would dominate an axis corner
or a pair corner; both are excluded by these bounds. \end{proof}

Let $E_{12}$ be the positive standard $xy$ points with $x+y>r$;
define $E_{13},E_{23}$ with thresholds $s,t$. Thus the positive plane
areas are $\binom r2+|E_{12}|$, and similarly for the other planes.

\begin{lemma}[disjoint resources]\label{res:A.4}
Put
\[
 C=\{(x,y,z)\in\mathbb Z_{>0}^3:x+y\le r,\ x+z\le s,\ y+z\le t\}.
\]
Without a full boundary, $C\subset D$. With full boundary
$g=(A,B,C_0)$, put
\[
 C^\circ=C\setminus(g+\mathbb N^3),\qquad
 L_+=\{q\in\mathbb Z_{>0}^3:q\le g,\ q\ne g\},\qquad H=C^\circ\cup L_+.
\]
Then $H\subset D$, and
\[
 m\ge a_1+a_2+a_3-2+
 \binom r2+\binom s2+\binom t2+
 |E_{12}|+|E_{13}|+|E_{23}|+|H|.\tag{3.1}\label{eq:capacity-1}
\]
Without a full boundary replace $H$ by $C$.
The proper lower box of $g$ consists of the vectors $q\le g$ other than
$g$; $L_+$ is its positive part. Removing the vectors above $g$ from $C$
gives the punctured core $C^\circ$.
\end{lemma}
\begin{proof}
All plane projections of a core point are standard by Lemma~\ref{res:A.3}. A boundary
below it must therefore have full support, hence must be $g$ if it
exists. This proves the core assertions. Every proper point below $g$
is standard. Finally count the disjoint support classes: origin and
axes, each plane's triangle and extras, and $H$. \end{proof}

We will use the exact row count
\[
 |C|=\sum_{\substack{y,z\ge1\\y+z\le t}}
 \min\{[r-y]_+,[s-z]_+\}.\tag{3.2}\label{eq:capacity-2}
\]
For example, at $r\ge s\ge t=4$ its six rows give
\[
 |C|=6s-\begin{cases}13&r=s,\\11&r=s+1,\\10&r\ge s+2.\end{cases}\tag{3.3}\label{eq:capacity-3}
\]

\begin{lemma}[terminal capacity]\label{res:A.5}
With full boundary $g=(A,B,C_0)$, one has $r,s,t\le A+B+C_0-2$.
Call an $xy$ corner $b$ \emph{high} when $b_x\ge A,b_y\ge B$.
Such a corner returns to $c_be_3$, where
\[
 a_3-C_0\le c_b<a_3.\tag{3.4}\label{eq:capacity-4}
\]
\end{lemma}
\begin{proof}
The return support follows from Lemma~\ref{res:A.1}. The standard point
$w=(b_x-A,b_y-B,0)$ has the residue of $(c_b+C_0)e_3$, since
$\sigma(g)=0$. If $c_b+C_0<a_3$, injectivity would equate these
distinct standard points. This proves \eqref{eq:capacity-4}, whose interval has length
$C_0$ because $C_0e_3\in D$.
Partition the $xy$ corners into $b_x<A$, then $b_x\ge A,b_y<B$,
then the remaining high corners. Their sizes are at most $A-1,B-1,C_0$:
use distinct first coordinates, second coordinates, and returns,
respectively. Permuting coordinates proves the other bounds.
\end{proof}

\begin{lemma}[separation of high corners]\label{res:A.6}
For two high $xy$ corners $b,c$ as in Lemma~\ref{res:A.5} with $b_x<c_x$,
one has $c_x-b_x>A$ or $b_y-c_y>B$.
\end{lemma}
\begin{proof}
Suppose $d=c_x-b_x\le A$ and $e=b_y-c_y\le B$; both are positive.
Their return heights differ by less than $C_0$. According to their
order, subtracting the residue identities equates the residues of
$(d,0,0),(0,e,c_c-c_b)$ or of $(d,0,c_b-c_c),(0,e,0)$.
Each pair consists of distinct proper points below $g$, contradicting
injectivity. \end{proof}

\begin{lemma}[plane gap formula]\label{res:A.7}
Adjoin $(X_0,Y_0)=(0,a_2)$ and $(X_{r+1},Y_{r+1})=(a_1,0)$.
For $0\le i\le r$ put
$u_i=X_{i+1}-X_i-1\ge0$, $v_i=Y_i-Y_{i+1}-1\ge0$. Then
\[
 |E_{12}|=\sum_{i=0}^r(r-i)u_i+\sum_{j=0}^r jv_j
             +\sum_{0\le i\le j\le r}u_iv_j.\tag{3.5}\label{eq:capacity-5}
\]
\end{lemma}
\begin{proof}
To count the plane, partition it into vertical strips between successive
corners. The standard plane consists of strips $X_i\le x<X_{i+1}$,
$0\le y<Y_i$. Its area is
$\sum_{i\le j}(1+u_i)(1+v_j)$.
Expand and subtract the axes, of size $a_1+a_2-1$, and the triangle,
of size $\binom r2$, using
$a_1=r+1+\sum u_i$, $a_2=r+1+\sum v_j$. \end{proof}

\begin{lemma}[plane excess bounds]\label{res:A.8}
For positive $A,B$, define
\[
\begin{split}
 Q(r;A,B)&=\#\{1\le x\le A,\ 1\le y\le B:x+y>r\},\\
 W(r;A,B)&=\sum_{i=A}^{[r-B]_+}\min\{A(r-i),Bi\},\\
 \Phi(r;A,B)&=\min_{0\le l\le r}\bigl((r-l)[A-l]_+
             +l[B+l-r]_++[A-l]_+[B+l-r]_+\bigr),\\
 F(r;A,B)&=\max\{1,Q(r;A,B),W(r;A,B),\Phi(r;A,B)\}.
\end{split}
\]
An empty interval contributes zero. The index $l$ in $\Phi$ records
a cut between successive plane corners; we call this a rank cut.
With full boundary $g=(A,B,C_0)$,
$|E_{12}|\ge F(r;A,B)$, and analogously for the other two planes.
\end{lemma}
\begin{proof}
For $A\le i\le r-B$, corners $i,i+1$ are high by Lemma~\ref{res:A.3}. By Lemma~\ref{res:A.6},
$u_i\ge A$ or $v_i\ge B$, so
$(r-i)u_i+iv_i\ge\min\{A(r-i),Bi\}$. Sum in \eqref{eq:capacity-5} to obtain $W$.
For $\Phi$, take the strip $X_l\le A<X_{l+1}$ containing the standard
point $(A,B,0)$. Set $U=\sum_{i\le l}u_i$, $V=\sum_{j\ge l}v_j$.
Then $U\ge[A-l]_+$ and $V\ge[B+l-r]_+$. Formula \eqref{eq:capacity-5} is at least
$(r-l)U+lV+UV$, yielding the stated minimum.
The whole positive rectangle below $(A,B,0)$ is standard, giving $Q$.
Finally, if $A+B\le r$, the $i=A$ term of $W$ is positive.
If $A+B>r$, every rank-cut cost is positive: for $l<A$ its first
term is positive unless $l=r$, when its cross term is positive;
for $l\ge A$ its second term is positive. Thus the unit bound also holds.
These bounds refer to the same excess, so we take their maximum.
\end{proof}

\begin{lemma}[joint axis gain]\label{res:A.9}
Order $r\ge s\ge t$ and put
$L_1=\max(r,A)$, $L_2=\max(r,B)$, $L_3=\max(r,C_0)$,
\[
\begin{split}
 h_{12}&=[r+A+B-L_1-L_2]_+,\\
 h_{13}&=[s+A+C_0-L_1-L_3]_+,\\
 h_{23}&=[t+B+C_0-L_2-L_3]_+,\\
 J&=\max\{h_{12},h_{13},h_{23},
               \lceil(h_{12}+h_{13}+h_{23})/2\rceil\}.
\end{split}
\]
The full resource bound, denoted by $\mathsf{BF}$, is
\[
\begin{split}
 m\ge\mathsf{BF}:={}&1+L_1+L_2+L_3+J+|H|
       +\binom r2+\binom s2+\binom t2\\
 &+F(r;A,B)+F(s;A,C_0)+F(t;B,C_0).
\end{split}\tag{3.6}\label{eq:capacity-6}
\]
\end{lemma}
\begin{proof}
The first two positive axis lengths $a_i-1$ are at least $r$ by Lemma~\ref{res:A.3};
the third is too, since the $r$ distinct $xy$ returns are positive
third-axis points (they cannot equal the zero return of $g$).
All axis lengths also dominate the corresponding coordinate of $g$.
For the standard face $(A,B,0)$, assign each $xy$ corner with first
coordinate $>A$ to that coordinate; all other corners have second
coordinate $>B$ and are assigned there. Hence
$r+A+B\le(a_1-1)+(a_2-1)$.
The other faces give the analogous inequalities. For
$\lambda_i=(a_i-1)-L_i\ge0$, these say $h_{ij}\le \lambda_i+\lambda_j$.
Thus $\sum \lambda_i\ge J$. Substitute into \eqref{eq:capacity-1} and apply Lemma~\ref{res:A.8}.
\end{proof}

\begin{lemma}[interior resources]\label{res:A.10}
One has
\[
 |H|\ge\max\{ABC_0-1,\ |C|-|C\cap(g+\mathbb N^3)|\}.
\]
A nonempty core hole $C\cap(g+\mathbb N^3)$ forces $r+4\le s+t$;
in particular the hole is empty when $t\le3$ and $r\ge s$.
\end{lemma}
\begin{proof}
The two entries count $L_+$ and $C^\circ$, respectively, both contained
in $H$. If a core point dominates $g$, summing its plane inequalities
gives $2(A+B+C_0)\le r+s+t$. Combine with $A+B+C_0\ge r+2$ from Lemma~\ref{res:A.5}.
\end{proof}

Let $r,s,t$ count the boundary points supported on the $xy,xz,yz$
planes. The coordinate transpositions permute these three counts, so
we may order them as $r\ge s\ge t$. Transport $D$, its residue weights
and its returns together. This preserves all hypotheses.
If a full-support boundary exists, write it as $g=(A,B,C_0)$.
By Lemma~\ref{res:A.2}, it suffices to establish the two capacities
\[
m\ge\left\lceil(r+s+t+2)^2/4\right\rceil
\quad\text{with a full boundary},\tag{3.7}\label{eq:capacity-7}
\]
\[
m\ge\left\lceil(r+s+t+1)^2/4\right\rceil
\quad\text{without one}.\tag{3.8}\label{eq:capacity-8}
\]
The following cases prove these inequalities using only the resources
above. The final square-root conversion is given at the end of this section.
\subsection{Staircases without a full-support boundary}
A plane is called \emph{exact} when its extra-point set $E_{ij}$ is empty,
so its positive standard points form exactly the forced triangle.
\begin{lemma}[exact planes and spacing of returns]\label{res:A.20} Call a pair family origin-free if none of its returns is zero. At least two of the three pair families are origin-free. If \(E_{12}=\varnothing\), its corners are exactly \((i+1,r-i,0)\), \(0\le i<r\). Write their distinct return coordinates as \(c_i\). An upward step has size at least \(t\); a downward step has size at least \(s\). For an exact \(xz\) plane the corresponding upward and downward lower bounds are \(t,r\); for an exact \(yz\) plane they are \(s,r\).
\end{lemma}
\begin{proof} Any two pair boundary points, including points from different pair families, share a positive coordinate. Lemma~\ref{res:A.1} therefore implies that at most one pair boundary point returns to zero. This proves the origin-free assertion. If there are no extras in the \(xy\) plane, its positive standard points are precisely \(x+y\le r\). Each \((i+1,r-i,0)\) is outside, while its positive-coordinate predecessors either lie in that triangle or on an axis at coordinate \(r\), which is standard by \(a_1,a_2\ge r+1\). Thus these are \(r\) boundary points and exhaust the pair family.

Subtracting consecutive corner residue identities gives
\(c_{i+1}p_3+p_2=c_ip_3+p_1\).
If \(d=c_{i+1}-c_i>0\), then \(p_1=p_2+dp_3\). If \(d<t\), the distinct points \(e_1\) and \((0,1,d)\) are both standard: the latter lies in the forced \(yz\) triangle and the former is standard because \(r\ge2\) whenever an adjacent step exists. Their residues agree, a contradiction. Hence \(d\ge t\). For a downward step of size \(d\), the identity becomes \(p_2=p_1+dp_3\); if \(d<s\), the standard points \(e_2\) and \((1,0,d)\) collide. Permuting coordinates gives the other two pairs of bounds. \end{proof}

\begin{lemma}[surplus forced by origin-free families]\label{res:A.21} Assume \(r\ge s\ge t\).
If \(r=s=R\) and \(t\ge2\), the combined axis and plane extras contribute at least two beyond the baseline axis lengths \((R,R,R)\).
If \(r>s\), use baseline positive axis lengths \((r,r,r-1)\), and define
\[
\delta_1=(a_1-1)-r,\quad\delta_2=(a_2-1)-r,
\quad\delta_3=(a_3-1)-(r-1),
\quad E=\delta_1+\delta_2+\delta_3+|E_{12}|+|E_{13}|+|E_{23}|.
\]
These quantities are nonnegative. An origin-free \(xy\) family forces \(\delta_3\ge1\). An origin-free \(xz\) family forces \(\delta_2+|E_{13}|\ge1\) whenever \(s\ge2\) and \(r-1<t(s-1)\). An origin-free \(yz\) family forces \(\delta_1+|E_{23}|\ge1\) whenever \(t\ge2\) and \(r-1<s(t-1)\).
\end{lemma}
\begin{proof} The first two axis baselines follow from Lemma~\ref{res:A.3}. The third follows because the \(r\) distinct \(xy\) returns lie among \(0,\ldots,a_3-1\), so \(a_3\ge r\). In the balanced case \(a_3\ge s+1=r+1\) already follows from the \(xz\) corners. An origin-free family uses distinct positive return positions; for \(xy\) this directly forces \(a_3-1\ge r\), proving the strict-case assertion.

For \(xz\), suppose both \(\delta_2\) and \(|E_{13}|\) vanish. Its exact return sequence has \(s\) distinct terms in \([1,r]\). By Lemma~\ref{res:A.20}, a downward step would have size at least \(r\), impossible in an interval of diameter \(r-1\). Thus every step is upward, of size at least \(t\), and the total span is at least \(t(s-1)>r-1\), also impossible. The \(yz\) assertion is identical with its stated bounds.

In the balanced case an origin-free \(xy\) family with no extra and no third-axis surplus has \(R\) positive returns in \([1,R]\). Downward steps of size \(R\) are impossible; upward steps of size at least \(t\ge2\) produce total span at least \(2(R-1)>R-1\), since \(R\ge2\). Hence its third-axis surplus plus \(|E_{12}|\) is at least one. The same reasoning treats \(xz\). For an origin-free exact \(yz\) family with zero first-axis surplus, even its first two distinct returns lie in \([1,R]\), while either direction requires a jump at least \(R\), impossible. Thus each origin-free family forces one unit in its own resource group (third axis plus \(xy\) extras, second axis plus \(xz\) extras, or first axis plus \(yz\) extras). These three groups are disjoint. Lemma~\ref{res:A.20} supplies two origin-free families, so two units can be added without overlap. \end{proof}

Define the simplex count
\[
\Pi(s,t)=\sum_{u=1}^{t-1}u(s-u)
=\frac{t(t-1)(3s-2t+1)}6\qquad(0\le t\le s).
\]
The empty cases \(t=0,1\) give zero. For \(t\ge2\), count points with \(y+z=u+1\le t\) and \(1\le x\le s-u\). There are \(u\) choices of \((y,z)\) and \(s-u\) of \(x\). These points belong to the core because \(x+y+z\le s+1\), so \(x+y\le s\le r\) and \(x+z\le s\). The polynomial formula follows from \(\sum_{u=1}^{n}u=n(n+1)/2\) and \(\sum_{u=1}^{n}u^2=n(n+1)(2n+1)/6\); both identities follow by induction on \(n\). In particular \(|C|\ge \Pi(s,t)\).

\begin{lemma}[no-full capacity]\label{res:A.22}\label{res:A.23}
If no full-support boundary exists and $r\ge s\ge t$, then
\[
 m\ge\left\lceil(r+s+t+1)^2/4\right\rceil.
\]
\end{lemma}
\begin{proof}
Put $n=r+s+t+1$ and
\[
 M_{\mathrm{nf}}=3r+\binom r2+\binom s2+\binom t2.
\]
The positive axis lengths are at least $r,r,r-1$: the last bound
follows from the $r$ distinct $xy$ returns on the third axis.
Thus the disjoint resources give $m\ge M_{\mathrm{nf}}+|C|\ge M_{\mathrm{nf}}+\Pi(s,t)$.

For every $t\ge5$, including $\delta=0$, set $\ell_{\Pi}=(t-1)s-t+5$.
With $\delta_{st}=s-t\ge0$ and $v=t-5\ge0$, direct expansion gives
\[
\begin{split}
 6\Pi(s,t)-6\ell_{\Pi}
 &=v(t^2-t+6)+3\delta_{st}(t-1)(t-2)\ge0,\\
 4(M_{\mathrm{nf}}+\ell_{\Pi}+1)-(n^2+3)&=(\delta-(t-4))^2+4>0.
\end{split}
\]
Since $M_{\mathrm{nf}}+\ell_{\Pi}$ is an integer and
$\lceil n^2/4\rceil=\lfloor(n^2+3)/4\rfloor$, these imply the target.

It remains that $t\le4$.  If $r=s=R$, the origin, axes and triangles
give $m\ge M=3R+1+2\binom R2+\binom t2$.
For $t=0,1$, compare this directly with $R^2+R+1,(R+1)^2$.
For $2\le t\le4$, Lemma~\ref{res:A.21} supplies two further axis/plane units,
and the core gives $\Pi(R,t)$, so
$m\ge N=3R+3+R(R-1)+\binom t2+\Pi(R,t)$.
The remaining differences are
\[
\begin{array}{c|c|c}
 t&\lceil(2R+t+1)^2/4\rceil&N-\lceil(2R+t+1)^2/4\rceil\\\hline
 2&R^2+3R+3&0\\
 3&R^2+4R+4&R-3\\
 4&R^2+5R+7&3R-12
\end{array}
\]
and are nonnegative because $R\ge t$.

Now let $r>s$, so $\delta\ge1$, and use the nonnegative surplus $E$ of
Lemma~\ref{res:A.21}.  Exact support accounting gives
\[
 m\ge M_{\mathrm{nf}}+|C|+E.\tag{3.9}\label{eq:capacity-9}
\]
The following lower bounds $|C|\ge H$ suffice in the indicated cases;
the last column is $4(M_{\mathrm{nf}}+H)-n^2$:
\[
\begin{array}{c|c|c}
 \text{condition}&H&\text{difference}\\\hline
 t=0&0&\delta^2+8\delta+4s-1\\
 t=1&0&\delta^2+6\delta-4\\
 t=2,\ \delta\ge2&s-1&\delta^2+4\delta-9\\
 t=3,\ \delta\ge2\text{ or }s\ge5&3s-4&\delta^2+2\delta+4s-20\\
 t=4,\ s\ge5&6s-14&\delta^2+12s-57
\end{array}
\]
All differences are nonnegative under their row conditions.
For $t=2,3$, the core row lengths are respectively $s-1$ and
$s-1,s-2,s-1$; the width-four bound is $\Pi(s,4)$.
When $s=t=4$, the three square slices have total area
$3^2+2^2+1^2=14$: the constraint $x+y\le r$ removes only
$(3,3,1)$ at $r=5$, and nothing for $r\ge6$.
Thus $M_{\mathrm{nf}}+|C|=50>49$ at $(5,4,4)$, while for $r\ge6$,
\[
 4(M_{\mathrm{nf}}+14)-(r+9)^2=(r-6)^2+4(r-6)+11\ge0.
\]

Only $t=2,\delta=1$ and $(r,s,t)=(4,3,3),(5,4,3)$ remain.
For $t=2,\delta=1$, the target is $M_{\mathrm{nf}}+s$ and $|C|\ge s-1$;
failure forces $E=0$.  The $xy$ family cannot be origin-free,
since it would force $\delta_3\ge1$, so the $xz$ family is origin-free.
For $s\ge3$, its span condition $r-1=s<2(s-1)$ contradicts $E=0$
by Lemma~\ref{res:A.21}.  At $s=2$, the exact $xy$ plane has three distinct
returns in $\{0,1,2\}$, whose consecutive differences have absolute
value at least two by Lemma~\ref{res:A.20}; no such three-term sequence exists.
At $(4,3,3)$ the target is $M_{\mathrm{nf}}+7$ and $|C|=5$, so failure gives
$E\le1$.  All three origin-free resource implications apply
($3<3\cdot2$); two origin-free families force two units in disjoint
resource groups, a contradiction.
At $(5,4,3)$ the target is $M_{\mathrm{nf}}+9$ and $|C|=8$, so failure gives $E=0$.
Again $xy$ cannot be origin-free, while the origin-free $xz$ family
has $4<3\cdot3$ and forces positive surplus.
This exhausts the cases. \end{proof}

\subsection{Full-support widths at most two}

Use the notation above, with $r\ge s\ge t$ and
$g=(A,B,C_0)$. Write $\Lambda=\lceil(r+s+t+2)^2/4\rceil$. The weakened resource inequality used below is
\[
\begin{split}
 m\ge{}&1+\max(r,A)+\max(r,B)+\max(r,C_0)+|H|\\
 &+\binom r2+\binom s2+\binom t2
   +F(r;A,B)+F(s;A,C_0)+F(t;B,C_0).
\end{split}\tag{3.10}\label{eq:capacity-10}
\]
The full bound of Lemma~\ref{res:A.9} additionally retains the joint axis gain $J$.

\begin{lemma}[full-support width zero]\label{res:A.19} If the counts are ordered and \(t=0\), \eqref{eq:capacity-7} holds.
\end{lemma}
\begin{proof} Lemma~\ref{res:A.9}, or just \eqref{eq:capacity-1} and the three full-support axis baselines, gives \(m\ge M=3r+1+\binom r2+\binom s2\). Expansion gives \(4M-(r+s+2)^2=\delta^2+6\delta\ge0\). Since \(M\) is an integer, \(\lceil(r+s+2)^2/4\rceil\le M\le m\). \end{proof}


\begin{lemma}[width one]\label{res:A12.1}
If $t=1$, then $m\ge \Lambda$.
\end{lemma}
\begin{proof}
For positive integers $A,B,C_0$, expansion about $(1,1,1)$ gives
$ABC_0\ge A+B+C_0-2$.  The proper box and terminal capacity therefore
give $|H|\ge ABC_0-1\ge r-1$.  The three axes contribute at least
$3r$, and the three plane excesses at least three, so
\[
 m\ge4r+3+\binom r2+\binom s2=:M,\qquad
 4M-(r+s+3)^2=\delta^2+8\delta+3\ge0.
\]
\end{proof}

For the rest of this subsection $t=2$, so $r\ge s\ge2$ and
$\Lambda=\lceil(r+s+4)^2/4\rceil$.

\begin{lemma}[the profile $(2,2,2)$ of pair counts]\label{res:A12.2}
If $r=s=t=2$, then $m\ge16=\Lambda$.
\end{lemma}
\begin{proof}
All intercepts are at least three, and no pair can both equal three:
their gap sums would vanish in Lemma~\ref{res:A.7}, making the plane excess zero,
contrary to Lemma~\ref{res:A.8}.  Thus each pair of intercepts has sum at least seven,
and the three positive axis lengths sum to at least eight.
The origin and axes contribute nine points; the three triangles and
three unit excesses contribute six.  Capacity gives
$A+B+C_0\ge4$, so $(1,1,1)$ is a positive proper point below $g$,
contributing the sixteenth point in $H$. \end{proof}

\begin{lemma}[no coordinate equal to one]\label{res:A12.3}
If $A,B,C_0\ge2$, then $m\ge \Lambda$.
\end{lemma}
\begin{proof}
Expansion about $(2,2,2)$ and capacity give
$ABC_0\ge4(A+B+C_0-4)\ge4(r-2)$.
Equation \eqref{eq:capacity-10} yields
\[
 m\ge3r+\binom r2+\binom s2+ABC_0+4=:M,
\]
and for $r\ge4$,
\[
 4M-(r+s+4)^2\ge\delta^2+10\delta+8r-32\ge0.
\]
For $r\le3$, the full proper box gives
$m\ge(A+1)(B+1)(C_0+1)-1\ge26>\Lambda$.
\end{proof}

\begin{lemma}[only the first coordinate equal to one]\label{res:A12.4}
If $A=1$, $B,C_0\ge2$, then $m\ge \Lambda$.
\end{lemma}
\begin{proof}
Both the $yz$ rectangle excess and the positive proper box contain
$BC_0-1$ points.  Keeping one excess on each other plane gives
\[
 m\ge3r+\binom r2+\binom s2+2BC_0+2=:M.
\]
Since $B+C_0\ge r+1$ and
$BC_0\ge2(B+C_0-2)\ge2r-2$, for $r\ge3$,
\[
 4M-(r+s+4)^2\ge\delta^2+10\delta+8r-24\ge0.
\]
The case $r=2$ is Lemma~\ref{res:A12.2}. \end{proof}

\begin{lemma}[one transverse coordinate equal to one]\label{res:A12.5}
If $A\ge2$ and exactly one of $B,C_0$ is one, then $m\ge \Lambda$.
\end{lemma}
\begin{proof}
Let $N\ge2$ be the other transverse coordinate, and put $n=r$ if
$B=1$, $n=s$ if $C_0=1$.  The proper box has $AN-1$ points,
the short-plane rectangle excess has $N-1$, and the plane with
dimensions $(A,1)$ contributes $D_n=F(n;A,1)$.  The remaining plane
contributes at least one.  Together with the axes and triangles,
\[
 m\ge1+3r+\binom r2+\binom s2+N(A+1)+D_n=:M.\tag{3.11}\label{eq:capacity-11}
\]
Take $r\ge3$, since Lemma~\ref{res:A12.2} covers $r=2$.
Capacity gives $A+N\ge r+1$, and
\[
 N(A+1)=2(A+N-1)+(N-2)(A-1)\ge2r.
\]
Thus $D_n\ge1$ gives
$4M-(r+s+4)^2\ge\delta^2+10\delta-8>0$ if $\delta\ge1$.

If $r=s$, then $n=r$, and $A<r$ gives $D_n\ge A$ from the $i=A$
weighted term.  We claim $N(A+1)+D_n\ge2r+3$, except at $r=3$,
$A=N=2$.  For $A\ge r$ this is immediate; for $A=r-1$ it follows
from $2r+(r-1)\ge2r+3$ when $r\ge4$, while $r=3,N\ge3$ gives
at least eleven.  For $A\le r-2$, $r\ge4$ and
\[
\begin{split}
 N(A+1)+D_n-(2r+3)
 &\ge(A-1)r-A^2+A-2\\
 &=(A-2)(r-A-1)+(r-4)\ge0.
\end{split}
\]
This proves the balanced target.  At the exception
$(r,s,t)=(3,3,2)$, $g=(2,1,2)$ or $(2,2,1)$, Lemma~\ref{res:A.9} has $J=1$,
the three triangle areas sum to seven, and $|H|\ge3$.
The count-three excesses at $(2,1)$ and $(2,2)$ are at least two
(respectively $W=2$ and $\Phi=2$), and the last excess is at least one.
Hence $m\ge1+9+1+7+3+2+2+1=26\ge25=\Lambda$.
\end{proof}

\begin{lemma}[two coordinates equal to one, including the first]\label{res:A12.6}
If $g=(1,1,N)$ or $(1,N,1)$, then $m\ge \Lambda$.
\end{lemma}
\begin{proof}
Again take $r\ge3$; capacity gives $N\ge r$.
For $g=(1,1,N)$, the axis maxima sum to $2r+N$,
$|H|\ge N-1$, and the $xz,yz$ rectangle extras are
$N+1-s,N-1$.  For $g=(1,N,1)$ the analogous extras are
$N+1-r,N-1$.  Using $4N-s\ge3r$ or $4N-r\ge3r$, respectively,
gives
\[
 m\ge1+5r+\binom r2+\binom s2+W+\chi_3,\tag{3.12}\label{eq:capacity-12}
\]
where $W=W(r;1,1)$ in the first orientation and $W(s;1,1)$ in the
second; $\chi_3=1$ when $r=s=3$, and zero otherwise.
The extra unit $\chi_3$ is the core point $(2,1,1)$, outside the $g$-box.
For $n\ge2$, $W(n;1,1)\ge1$; it equals two at $n=3$ and is at least
three for $n\ge4$.  Thus
\[
 4\left(1+5r+\binom r2+\binom s2+W+\chi_3\right)-(r+s+4)^2
 =\delta^2+10\delta-12+4(W+\chi_3)\ge0:
\]
if $\delta\ge1$ the right side is at least three, and if $\delta=0$
one has $W+\chi_3\ge3$. \end{proof}

\begin{lemma}[the full-support point $(k,1,1)$]\label{res:A12.7}
If $g=(k,1,1)$, then $m\ge \Lambda$.
\end{lemma}
\begin{proof}
Capacity gives $k\ge r\ge s\ge2$.  Write $\ell_i=a_i-1$ and
$e_{ij}=|E_{ij}|$.  The proper box and its two face rows give
\[
 |H|\ge k-1,\qquad e_{12}\ge k-r+1,\qquad e_{13}\ge k-s+1.
\]
There is a joint extra unit
\[
 (\ell_1-k)+\bigl(e_{12}-(k-r+1)\bigr)\ge1.
\]
Indeed, either $\ell_1>k$, or $\ell_1=k$ and the standard point
$(k,1,0)$ forces the last $xy$ corner to have height at least two.
Then $v_r\ge1$ in Lemma~\ref{res:A.7}, whose linear and cross terms give
$e_{12}\ge r+\sum u_i=k\ge(k-r+1)+1$.

Since $e_{23}\ge1$, at least one of $(0,1,2),(0,2,1)$ is standard.
Lift such a base to $1\le x\le s-2$.  Both long-plane projections
are in their forced triangles, and $x<k$ excludes $g$, so all $s-2$
points are standard.  They lie outside $H$: their transverse sum is
three, exceeding $t=2$, and one transverse coordinate exceeds the
$g$-box bound.  These columns are disjoint from all axis and plane
resources.  With $\ell_1\ge k$, $\ell_2,\ell_3\ge r$, the joint unit
and $e_{23}\ge1$, accounting gives
\[
 m\ge\binom r2+\binom s2+4k+r+3=:M,\qquad
 4M-(r+s+4)^2=\delta^2+10\delta+16(k-r)-4.\tag{3.13}\label{eq:capacity-13}
\]
This suffices unless $r=s=k$.  The case $k=2$ is Lemma~\ref{res:A12.2}.
For $k\ge3$, $M=(k+1)(k+3)=\Lambda-1$, so failure forces $m=M$.
If $\ell_1=k$, the last-gap estimate above strengthens the joint
payment from one to $k-1\ge2$, impossible at this value of $m$.
Thus $\ell_1\ge k+1$.  The disjoint resources consisting of the
origin; three axes; three positive planes; $H$; and the chosen column
have lower bounds
\[
 \left(1;\ k+1,k,k;\ \binom k2+1,\binom k2+1,2;\ k-1;\ k-2\right).
\]
Their sum is $m$; hence all nonnegative slacks vanish.  In particular
\[
 (\ell_1,\ell_2,\ell_3)=(k+1,k,k),\qquad
 e_{12}=e_{13}=e_{23}=1.
\]
The unique $xy$ extra is $(k,1,0)$, so $(k+1,1,0)$ is a boundary
point and returns to $ce_3$ with $0\le c\le k$.  Its residue identity
and the terminal relation give
\[
 cp_3=(k+1)p_1+p_2,\quad kp_1+p_2+p_3=0,
 \qquad p_1=(c+1)p_3.
\]
Injectivity at $e_1$ and the pure $z$-axis points forces $c=k$, so
$p_1=(k+1)p_3$.
The unique $yz$ extra is one of $(0,1,2),(0,2,1)$; the other is a
boundary point, returning to $ae_1$.  Thus either
$2p_2+p_3=ap_1$ or $p_2+2p_3=ap_1$.
Reduce these equations modulo $k+1$, which divides
$m=(k+1)(k+3)$.  First $p_1=0$, then the terminal relation gives
$p_2+p_3=0$, and either $yz$ relation forces $p_2=p_3=0$.
All standard-point residues would be zero modulo $k+1>1$,
contradicting surjectivity.  This excludes failure. \end{proof}

\begin{theorem}[width-two capacity]\label{res:A12.8}
Every ordered full-support cyclic staircase with $r\ge s\ge t=2$
satisfies $m\ge\lceil(r+s+4)^2/4\rceil$.
\end{theorem}
\begin{proof}
If no coordinate of $g$ is one, use Lemma~\ref{res:A12.3}.  If $A=1$, use Lemma~\ref{res:A12.4} or
Lemma~\ref{res:A12.6} according as neither or at least one transverse coordinate is
one.  If $A\ge2$, use Lemma~\ref{res:A12.5} or Lemma~\ref{res:A12.7} according as exactly one or both
transverse coordinates are one.  Capacity excludes $(1,1,1)$.
\end{proof}

\subsection{Full-support width three}

Throughout this subsection $r\ge s\ge t=3$, $g=(A,B,C_0)$,
$N=\binom r2+\binom s2$, and $\mathcal Q=(r+s+5)^2$. Write $\Lambda=\lceil\mathcal Q/4\rceil$ for the integer target.
We prove $m\ge\lceil\mathcal Q/4\rceil$ from the resources above.
Use $E_n(U,V)=Q(n;U,V)$ for the rectangle excess of Lemma~\ref{res:A.8}, and put
\[
 d_n(U,V)=\max\{1,E_n(U,V),W(n;U,V)\}.
\]
For width three we also use the reduced resource bound $\mathsf{BW}$.
Omitting $J$ and the rank-cut improvement from Lemma~\ref{res:A.9} gives
\[
\begin{split}
 m\ge\mathsf{BW}:={}&1+\max(r,A)+\max(r,B)+\max(r,C_0)+N+3+|H|\\
 &+d_r(A,B)+d_s(A,C_0)+d_3(B,C_0).
\end{split}\tag{3.14}\label{eq:capacity-14}
\]
The full resource $\mathsf{BF}$ retains $J$ and replaces $d$ by $F$.
Terminal capacity is $A+B+C_0\ge r+2$; in particular $g\ne(1,1,1)$.

\begin{lemma}[short-plane resources]\label{res:A3.1}
For $U,V\ge2$ one has $E_3(U,V)=UV-3$.
If $U+j\le n$ for $j=1,2$, then $W(n;U,1)\ge jU$.
Also $d_3(U,1)\ge[U-2]_+$, $d_3(2,1)=2$, and $d_3(3,1)=1$.
If $U,V\ge2$ and $U+V\ge n+2$, then $E_n(U,V)\ge3$;
if $U+V=n+1$, then $E_n=1$, $W=0$, $d_n=1$, and
$\Phi(n;U,V)=\min(U,V)$. We call this last case the first diagonal.
With $\ell_{\mathrm{row}}=\min(r-2,s-1)$, the following counts are exact:
\[
\begin{array}{c|l}
g&|H|\\\hline
(A,B,1),\ A,B\ge2&AB+s-3+[s-1-A]_++[\ell_{\mathrm{row}}-A]_+\\
(A,1,C_0),\ A,C_0\ge2&AC_0-1+\ell_{\mathrm{row}}+[s-1-A]_++[s-2-A]_+\\
(1,B,C_0),\ B,C_0\ge2&BC_0-1+(s-2)+(s-3)+(\ell_{\mathrm{row}}-1)
\end{array}\tag{3.15}\label{eq:capacity-15}
\]
\end{lemma}
\begin{proof}
The short triangle consists of $(1,1),(1,2),(2,1)$; the first two
weighted terms prove the $W$ bounds.  The three rectangle points
$(U,V),(U-1,V),(U,V-1)$ prove the three-point bound, and substitution
in the rank-cut formula proves the identity on the first diagonal.
For \eqref{eq:capacity-15}, the core has exactly the three rows $(y,z)=(1,1),(1,2),(2,1)$,
with $x$-endpoints $s-1,s-2,\ell_{\mathrm{row}}$.  Add their portions outside the positive
proper $g$-box, whose size is $ABC_0-1$.  This gives the displayed tails;
the rows are disjoint.  In the second row, capacity
$A+C_0\ge r+1$ ensures that a tail at $z=2$ is not removed by the puncture
when $C_0=2$. \end{proof}

\begin{lemma}[no coordinate equal to one]\label{res:A3.2}
If $A,B,C_0\ge2$, then $m\ge\lceil \mathcal Q/4\rceil$.
\end{lemma}
\begin{proof}
Set $V_{\mathrm{box}}=(A+1)BC_0$.  The positive box and short-plane rectangle in
\eqref{eq:capacity-14} give
\[
 m\ge N+3r+V_{\mathrm{box}}+d_r(A,B)+d_s(A,C_0).\tag{3.16}\label{eq:capacity-16}
\]
Consequently $V_{\mathrm{box}}\ge3r+5$ suffices, since
\[
 4(N+3r+V_{\mathrm{box}}+2)-\mathcal Q=\delta^2+4V_{\mathrm{box}}-12s-17\ge \delta^2+12\delta+3.\tag{3.17}\label{eq:capacity-17}
\]
For $r\ge6$ this condition follows from the estimate
\[
\begin{array}{c|l}
A\ge r-1&V_{\mathrm{box}}\ge4r\ge3r+5\\
2\le A\le r-3&BC_0\ge2(r-A),\quad
 (r-A)(A+1)-3(r-2)=(A-2)(r-3-A)\ge0,\\
 &\hfill V_{\mathrm{box}}\ge6(r-2)\ge3r+5\\
A=r-2,\ BC_0\ge6&V_{\mathrm{box}}\ge6(r-1)\ge3r+5.
\end{array}
\]
In the remaining case $A=r-2$, $B=C_0=2$, the term $i=r-2$ gives
$d_r\ge2r-4$, improving its baseline one by $2r-5$; now
$V_{\mathrm{box}}+(2r-5)=6r-9\ge3r+5$.

For $r=5$, one has $V_{\mathrm{box}}+d_5(A,B)\ge21$: if $A\ge4$ this is immediate;
if $A=3$, either $BC_0\ge6$ or $(B,C_0)=(2,2)$, where $V_{\mathrm{box}}+W_5=16+6$;
if $A=2$, either $BC_0\ge8$ or capacity leaves $(B,C_0)=(2,3),(3,2)$,
giving $V_{\mathrm{box}}+d_5\ge26,24$.  Hence $m\ge\binom s2+47$, giving
$50,53,57$ for $s=3,4,5$, against targets $43,49,57$.
For $r=3$, failure in \eqref{eq:capacity-16} forces $V_{\mathrm{box}}\le13$, hence $g=(2,2,2)$.
For $r=4$, failure forces $A\le3$, $BC_0\le7$; the complete remaining
table from \eqref{eq:capacity-16} is
\[
\begin{array}{c|rr}
g&s=3\ (\Lambda=36)&s=4\ (\Lambda=43)\\\hline
(2,2,2)&38&44\\
(2,2,3)&46&47\\
(2,3,2)&41&47\\
(3,2,2)&41&42\\
(3,2,3)&52&52\\
(3,3,2)&51&52
\end{array}
\]
The two residual cases are treated next. \end{proof}

\begin{lemma}[the two residual cases with all coordinates at least two]\label{res:A3.3}
The residual cases $r=s=3$, $g=(2,2,2)$ and $r=s=4$, $g=(3,2,2)$
satisfy their targets.
\end{lemma}
\begin{proof}
For $g=(2,2,2)$, the proper box $[0,2]^3\setminus\{g\}$ and the three
axis points at height three give $29$ standard points, so only $m=29,30$
could fail.  Write $a,b,c$ for the weights.  The relation
$2(a+b+c)=0$ forces $a+b+c=0$ if $m=29$, contradicting injectivity at
$(1,1,1)$ and zero.  For $m=30$, the only alternative is $a+b+c=15$.
If exactly one weight is odd the displayed set has $19$ even residues;
if all three are odd it has $16$ odd residues.  Either exceeds $15$.

For $g=(3,2,2)$, its $35$-point proper box and the seven standard points
\[
(4,0,0),(0,3,0),(0,4,0),(0,0,3),(0,0,4),(1,3,0),(1,0,3)
\]
give $m\ge42$, so failure requires $m=42$.  Their membership follows
from the height-four axes and long-plane triangles.  The terminal
relation $3a+2b+2c=0$ makes $a$ even.  If $b,c$ are both even, all
$42$ residues are even; if exactly one is odd, the proper box alone has
$23$ even residues; if both are odd, it has $19$, and the three added
height-four axis points raise this to $22$.  Each exceeds the $21$
even residue classes. \end{proof}

\begin{lemma}[one transverse coordinate equal to one]\label{res:A3.4}
If $C_0=1$, $A,B\ge2$, or $B=1$, $A,C_0\ge2$, then
$m\ge\lceil \mathcal Q/4\rceil$.
\end{lemma}
\begin{proof}
First let $C_0=1$ and put
$\xi_{\mathrm{rect}}=E_r(A,B)$, $\xi_{\mathrm{long}}=d_s(A,1)$, $\xi_{\mathrm{short}}=d_3(B,1)$,
$\xi_{\mathrm{end}}=[s-1-A]_+$, $\xi_{\mathrm{row}}=[\ell_{\mathrm{row}}-A]_+$.  Equations \eqref{eq:capacity-14}--\eqref{eq:capacity-15} give
\[
 m\ge1+3r+N+AB+s+\xi_{\mathrm{end}}+\xi_{\mathrm{row}}+\xi_{\mathrm{rect}}+\xi_{\mathrm{long}}+\xi_{\mathrm{short}}.\tag{3.18}\label{eq:capacity-18}
\]
If $A+B\ge r+2$, then $AB\ge2r$, $\xi_{\mathrm{rect}}\ge3$, and $\xi_{\mathrm{long}},\xi_{\mathrm{short}}\ge1$, whence
$4m-\mathcal Q\ge \delta^2+8\delta+4(\xi_{\mathrm{rect}}+\xi_{\mathrm{long}}+\xi_{\mathrm{short}})-21>0$ for $\delta\ge1$.
For $\delta=0$ one has $\xi_{\mathrm{rect}}+\xi_{\mathrm{long}}+\xi_{\mathrm{short}}\ge6$: $B=2$ gives $\xi_{\mathrm{short}}=2$;
$B=3$ gives either $\xi_{\mathrm{long}}\ge A\ge2$ or, when $A\ge r$, four rectangle
points $(A,1),(A,2),(A,3),(A-1,3)$; $B\ge4$ gives $\xi_{\mathrm{short}}\ge B-2\ge2$.

It remains that $A+B=r+1$.  Put $Z=\xi_{\mathrm{end}}+\xi_{\mathrm{row}}+\xi_{\mathrm{rect}}+\xi_{\mathrm{long}}+\xi_{\mathrm{short}}$; then
\[
 4m-\mathcal Q\ge \delta^2+8\delta+4Z-29.\tag{3.19}\label{eq:capacity-19}
\]
For $\delta\ge2$, $Z\ge3$ suffices.  The bounds needed at $\delta=0,1$ are
\[
\begin{array}{c|c|c}
\delta&\text{placement}&\text{bound}\\\hline
0&B=2&Z\ge r+2\quad(r\ge6)\\
0&B=3&Z\ge2r-1\quad(r\ge5)\\
0&B\ge4&Z\ge r+2B-5\ge8\\
1&B\ge3&Z\ge r+2B-6\ge5\quad(r\ge5).
\end{array}
\]
Indeed, the first two rows use $(\xi_{\mathrm{end}},\xi_{\mathrm{row}},\xi_{\mathrm{long}},\xi_{\mathrm{short}})=(0,0,\ge A,2)$ and
$(1,0,\ge2A,\ge1)$; the next uses
$(B-2,B-3,\ge A,\ge B-2)$, and the last uses
$(B-3,B-3,\ge A,\ge B-2)$, always with $\xi_{\mathrm{rect}}\ge1$.

The remaining family is $s=A=r-1$, $B=2$, $C_0=1$, $r\ge4$.
Its disjoint $H$ rows have lengths $r-1,r-3,r-2$ at bases
$(y,z)=(1,1),(1,2),(2,1)$.  Also $\Phi_r(A,2)=2$,
$d_3(2,1)=2$, the remaining excess is at least one, and $J=1$
($(h_{12},h_{13},h_{23})=(1,0,0)$).  Therefore
\[
 m\ge\binom r2+\binom{r-1}2+6r+4=(r+2)^2+1>\mathcal Q/4.
\]
The other residual placements are the following finite table:
\[
\begin{array}{c|r|r|r}
(r,s);g&\mathsf{BW}&\mathsf{BF}\text{ lower bound}&\lceil \mathcal Q/4\rceil\\\hline
(4,3);(2,3,1)&35&37&36\\
(4,4);(3,2,1)&41&43&43\\
(5,5);(4,2,1)&56&58&57\\
(4,4);(2,3,1)&42&44&43\\
(3,3);(2,2,1)&28&30&31
\end{array}
\]
Here the two added units are the first-diagonal rank-cut improvement
and $J\ge1$.  Lemma~\ref{res:A3.8} excludes $m=30$ in the final row.

For $B=1$, put $\xi_{\mathrm{rect}}=E_s(A,C_0)$, $\xi_{\mathrm{long}}=d_r(A,1)$, $\xi_{\mathrm{short}}=d_3(1,C_0)$.
The corresponding resource is
\[
\begin{split}
 m\ge{}&3r+N+AC_0+\ell_{\mathrm{row}}+3+[s-1-A]_++[s-2-A]_+\\
 &+\xi_{\mathrm{long}}+\xi_{\mathrm{rect}}+\xi_{\mathrm{short}}.
\end{split}\tag{3.20}\label{eq:capacity-20}
\]
When $r=s$ this is \eqref{eq:capacity-18} with $B$ replaced by $C_0$.
The slack case $A+C_0\ge r+2$ follows from the preceding slack estimate
using $\ell_{\mathrm{row}}\ge s-2$.  In the remaining tight case $A+C_0=r+1$, $r>s$,
one has $\ell_{\mathrm{row}}=s-1$, $\xi_{\mathrm{rect}}\ge3$, $\xi_{\mathrm{long}}\ge A$, $\xi_{\mathrm{short}}\ge C_0-2$, and
$AC_0\ge2r-2$.  Thus
\[
 m\ge N+6r+s+1,\qquad
 4(N+6r+s+1)-\mathcal Q=\delta^2+12\delta+4s-21\ge4.
\]
The balanced finite placements use the same table with $y,z$ exchanged.
\end{proof}

\begin{lemma}[only the first coordinate equal to one]\label{res:A3.5}
If $A=1$, $B,C_0\ge2$, then $m\ge\lceil \mathcal Q/4\rceil$.
\end{lemma}
\begin{proof}
The last row of \eqref{eq:capacity-15}, $\ell_{\mathrm{row}}\ge s-2$, and $E_3=BC_0-3$ give
$m+6\ge N+3r+2BC_0+3s$.  Capacity implies
$BC_0\ge2r-2$, so
\[
 m\ge N+7r+3s-10,\qquad
 4(N+7r+3s-10)-\mathcal Q=\delta^2+16r-65.
\]
This suffices for $r\ge5$ and $(r,s)=(4,3)$.
At $(r,s)=(3,3),(4,4)$, $BC_0\ge7$ already gives $m\ge32,44$.
Otherwise $(B,C_0)=(2,2),(2,3),(3,2)$, with $(2,2)$ excluded by
capacity when $r=4$.  The latter two placements give
$\mathsf{BW}=31$ at $r=s=3$ and $47$ at $r=s=4$.
Only $r=s=3$, $g=(1,2,2)$ remains: $\mathsf{BW}=28$,
$\Phi_3(2,2)=2$ improves its unit excess by one, and $J=1$ gives
$\mathsf{BF}\ge30$; Lemma~\ref{res:A3.8} excludes $m=30$. \end{proof}

\begin{lemma}[the full-support point $(h,1,1)$]\label{res:A3.6}
If $g=(h,1,1)$, then $m\ge\lceil \mathcal Q/4\rceil$.
\end{lemma}
\begin{proof}
Capacity gives $h\ge r$.  If $r>s$, the three $H$ rows have lengths
$h-1,s-2,s-1$.  With $E_r(h,1)=h+1-r$, $E_s(h,1)=h+1-s$,
and $d_3(1,1)=2$, \eqref{eq:capacity-14} yields
\[
 m\ge N+r+s+4h+4,\qquad
 4(N+r+s+4h+4)-\mathcal Q=(\delta-1)(\delta+9)+16(h-r)\ge0.
\]
Now let $r=s=k\ge4$.  The row lengths are $h-1,k-2,k-2$, giving
$m\ge k^2+k+4h+3$, sufficient if $h\ge k+1$.
For $h=k$, write $A_{yz}$ for the number of standard points with $y,z>0$ in the $yz$ plane.
Each standard base $(0,u,v)$ with $u+v>3$ lifts to the disjoint column
\[
 (x,u,v),\qquad 1\le x\le[k-\max(u,v)]_+.
\]
Its long-plane projections lie in their forced triangles, and $x<k$
excludes $g$; hence it is standard.  These columns lie outside both the
core and the $g$-box.  For any finite collection $B_0$ of such bases,
\[
 m\ge k^2+5k-2+A_{yz}
       +\sum_{(0,u,v)\in B_0}[k-\max(u,v)]_+.\tag{3.21}\label{eq:capacity-21}
\]
The short excess gives $A_{yz}\ge5$.  If both $(0,1,3),(0,3,1)$ are
standard, their columns give $m\ge k^2+7k-3\ge k^2+5k+7$ for $k\ge5$.
If $(0,1,3)$ is missing, the three short-plane corner heights must be
$3,2,1$.  The two internal gaps then have zero vertical excess, so Lemma~\ref{res:A.8}
forces horizontal excess at least one in each.  Their predecessors
give the three extra bases $(0,2,2),(0,3,1),(0,4,1)$.
Thus $A_{yz}\ge6$, and their column lengths sum to $3k-9$ for $k\ge4$;
\eqref{eq:capacity-21} gives $m\ge k^2+8k-5\ge k^2+5k+7$.
Missing $(0,3,1)$ is treated by reflection.

Only $k=h=4$ with both bases standard remains.  Suppose $m\le42$,
and let $X,Y,Z$ be the intercepts.  All are at least five.
If $X=5$, the standard point $(4,1,0)$ forces the last $xy$ corner to
have height at least two, so its final gap contributes at least four
to that plane's excess.  In the basic resource sum this raises the
bound from $41$ to at least $44$, impossible.  Hence $X\ge6$.
The disjoint origin, axes, positive planes, $H$, and two columns now
have respective lower bounds
\[
 (1;\ 5,4,4;\ 7,7,5;\ 7;\ 2),
\]
summing to $42$.  All are attained, so $m=42$,
$(X,Y,Z)=(6,5,5)$ and $E_{xy}=E_{xz}=1$.
The unique long-plane extras are $(4,1,0),(4,0,1)$, making
$(5,1,0),(5,0,1)$ boundary points.  If their opposite-axis returns
have heights $c_0,d_0\le4$, then, writing $a,b,c$ for the weights,
\[
 4a+b+c=0,\quad c_0c=5a+b,\quad d_0b=5a+c,
 \quad a=(c_0+1)c=(d_0+1)b.
\]
Injectivity at $e_x$ and the pure-axis points forces $c_0=d_0=4$.
Adding the return identities gives $22a=0$ modulo $42$, hence $2a=0$,
contradicting injectivity at $2e_x$ and zero.

Finally, at $r=s=3$ the three $H$ rows give $|H|=h+1$, and
$\mathsf{BW}=2h+19+2\max(1,h-2)$.
This is at least $31$ for $h\ge4$; for $h=3$ it is $27$, and Lemma~\ref{res:A3.8}
excludes $m=27,28,29,30$. \end{proof}

\begin{lemma}[the full-support points $(1,h,1)$ and $(1,1,h)$]\label{res:A3.7}
Both forms of the full-support point satisfy $m\ge\lceil \mathcal Q/4\rceil$.
\end{lemma}
\begin{proof}
Here $h\ge r$, and direct counting of the same three rows gives
\[
\begin{array}{ll}
g=(1,h,1):&\mathsf{BW}+3=N+r+2s+\ell_{\mathrm{row}}+4h+W(s;1,1),\\
g=(1,1,h):&\mathsf{BW}+3=N+2r+s+\ell_{\mathrm{row}}+4h+W(r;1,1).
\end{array}\tag{3.22}\label{eq:capacity-22}
\]
Write $w$ for the respective $W$ term.  Its indices $1,n-1$ give
$w\ge2$ for $n\ge3$, and the index $2$ gives $w\ge4$ for $n\ge4$.
Since $\ell_{\mathrm{row}}\ge s-2$,
\[
 m\ge N+5r+3s+w-5,\qquad
 4(N+5r+3s+w-5)-\mathcal Q=\delta^2+8r+4w-45.
\]
This is positive for $r\ge5$.  If $r>s$, $\ell_{\mathrm{row}}=s-1$ improves the
constant $-5$ to $-4$; this gives $m\ge36$ at $(r,s)=(4,3)$.
At $(4,4)$, $w\ge4$ gives $m\ge43$.
At $(3,3)$ the exact identity is $\mathsf{BW}=4h+15$, sufficient for
$h\ge4$; Lemma~\ref{res:A3.8} excludes the two placements with $h=3$ and $m=27,\ldots,30$.
\end{proof}

\begin{lemma}[finite residue obstructions]\label{res:A3.8}
If $r=s=t=3$, a full boundary which is a permutation of $(2,2,1)$
cannot occur at $m=30$, and a permutation of $(3,1,1)$ cannot occur at
$m\in\{27,28,29,30\}$.
\end{lemma}
\begin{proof}
Let
\[
 \mathcal K=\{q\in\{0,1,2,3\}^3:|\operatorname{supp}(q)|\le1\}
 \cup\{q\in\mathbb N^3:|\operatorname{supp}(q)|=2,\ |q|_1\le3\},
\]
and define the following sets of $24$ and $25$ points, respectively:
\[
\begin{aligned}
 \mathcal V={}&\mathcal K\cup\{(2,2,0),(1,1,1),(1,1,2),(1,2,1),(2,1,1)\},\\
 \mathcal W={}&\mathcal K\cup\{(1,1,1),(2,1,1),(1,2,1),(1,1,2),(3,1,0),(3,0,1)\}.
\end{aligned}
\]
All points of $\mathcal K$ are standard by the axis bounds and forced triangles.
For $g=(2,2,1)$ the additional $\mathcal V$ points lie in its proper lower box,
except $(1,1,2)$, which lies in the punctured core.  For $g=(3,1,1)$
the additional $\mathcal W$ points lie in its proper lower box or punctured core.
Thus $\mathcal V\subset D$ or $\mathcal W\subset D$, as appropriate.  Permuting coordinates
handles the other placements, since the three pair counts are equal.

Write $a,b,c$ for the correspondingly relabelled residue weights.
The terminal relation gives $c=-2a-2b$ for $\mathcal V$, or $c=-3a-b$ for $\mathcal W$.
For a finite set $\mathcal U$ put
\[
 N_{\mathcal U,m}^{\alpha,\beta}(a,b)
 =\#\{(x-\alpha z)a+(y-\beta z)b\pmod m:(x,y,z)\in \mathcal U\}.
\]
The following complete finite arithmetic table gives the maximum over
\emph{all} $0\le a,b<m$:
\[
\begin{array}{c|c|c|c|c}
 \mathcal U&(\alpha,\beta)&m&\#\{(a,b)\}&\max N_{\mathcal U,m}^{\alpha,\beta}(a,b)\\\hline
 \mathcal V&(2,2)&30&900&23\\
 \mathcal W&(3,1)&27&729&23\\
 \mathcal W&(3,1)&28&784&24\\
 \mathcal W&(3,1)&29&841&23\\
 \mathcal W&(3,1)&30&900&23
\end{array}
\]
For completeness, an exact algorithm to reproduce each last entry is
\[
\begin{array}{l}
 N_{\max}\gets0;\\
 \textbf{for }a,b=0,\ldots,m-1:\\
 \quad \mathcal R_{\mathrm{img}}\gets\varnothing;\\
 \quad\textbf{for }(x,y,z)\in \mathcal U:\
       \mathcal R_{\mathrm{img}}\gets \mathcal R_{\mathrm{img}}\cup\{[(x-\alpha z)a+(y-\beta z)b]_m\};\\
 \quad N_{\max}\gets\max(N_{\max},\#\mathcal R_{\mathrm{img}});\\
 \textbf{return }N_{\max},
\end{array}
\]
where $[u]_m$ is the least nonnegative remainder.  The supplementary finite-residue verifier performs these $4154$
weight-pair checks using integer arithmetic.  Each maximum is smaller than $|\mathcal U|$, so two
distinct standard points have equal residue, contradicting injectivity.
\end{proof}


\begin{theorem}[width-three capacity]\label{res:A3.9}
Every ordered full-support cyclic staircase with $r\ge s\ge t=3$ satisfies
$m\ge\lceil(r+s+5)^2/4\rceil$.
\end{theorem}
\begin{proof}
Capacity excludes $(1,1,1)$.  According as zero, one, or two coordinates
of $g$ equal one, apply Lemma~\ref{res:A3.2}--Lemma~\ref{res:A3.3}, Lemma~\ref{res:A3.4}--Lemma~\ref{res:A3.5}, or Lemma~\ref{res:A3.6}--Lemma~\ref{res:A3.7}.
Lemma~\ref{res:A3.8} supplies every residual finite obstruction. \end{proof}


\subsection{Full-support widths at least four}
\begin{lemma}[quantitative core-hole capacity]\label{res:A.30} In the ordered full-support case put
\[
h=\left\lfloor\frac{[s+t+2-r]_+}{2}\right\rfloor,
\qquad c(t)=\binom{\lfloor(t+2)/2\rfloor}{3}.
\]
Then the core hole has cardinality at most \(\binom h3\le c(t)\). Consequently
\[
|H|\ge\max\{ABC_0-1,[|C|-c(t)]_+\},\qquad
m\ge M+\Pi(s,t)-c(t), \tag{3.23}\label{eq:capacity-23}
\]
where \(M=3r+1+\binom r2+\binom s2+\binom t2\).
\end{lemma}
\begin{proof} Translate a hole point \(q\ge g\) to the positive triple \((X,Y,Z)=q-g+(1,1,1)\). This map is injective. Summing the core inequalities gives \(2(q_x+q_y+q_z)\le r+s+t\), while Lemma~\ref{res:A.5} gives \(A+B+C_0\ge r+2\). Therefore \(2(X+Y+Z)\le s+t+2-r\), and \(X+Y+Z\le h\). The number of positive triples with sum at most \(h\) is \(\binom h3\): replace the triple by the three strictly increasing partial sums \(X,X+Y,X+Y+Z\) in \(\{1,\ldots,h\}\). This is a bijection with three-element subsets. Since \(s\le r\), \(h\le\lfloor(t+2)/2\rfloor\), proving the bound, including cases where the numerator truncates to zero. Lemma~\ref{res:A.10} gives the maximum. Counting only the punctured simplex inside the punctured core loses at most the same number of hole points; Lemma~\ref{res:A.4} and the full-support axis baselines yield the last inequality. \end{proof}

\begin{lemma}[all full-support widths at least seven]\label{res:A.31} For \(r\ge s\ge t\ge7\), \eqref{eq:capacity-7} holds.
\end{lemma}
\begin{proof} Put \(\Lambda=\lceil(r+s+t+2)^2/4\rceil\) and \(e=(r+s+t+2)\bmod2\). Expanding the square and using \(4\lceil n^2/4\rceil=n^2+3(n\bmod2)\) gives
\[
4\Lambda+(r-s-t+3)^2=4M+4st+9+3e.
\]
Thus \(\Lambda\le M+st+3\). It suffices by \eqref{eq:capacity-23} to prove \(\Pi(s,t)\ge st+3+c(t)\). For \(t=7\), \(\Pi(s,7)=21s-91\), \(c(7)=4\), and the difference is \(14s-98\ge0\).

For \(t\ge8\), put \(\delta_{st}=s-t\ge0\). If \(t=2n\), \(n\ge4\), then
\[
6\big(\Pi(s,t)-st-3-c(t)\big)
=7n^3-24n^2-n-18+6n(2n-3)\delta_{st}.
\]
Write \(n=4+k\), \(k\ge0\). Its cubic part becomes \(7k^3+60k^2+143k+42>0\). If \(t=2n+1\), \(n\ge4\), the corresponding identity is
\[
6\big(\Pi(s,t)-st-3-c(t)\big)
=7n^3-12n^2-19n-24+3(2n+1)(2n-2)\delta_{st}.
\]
The cubic part becomes \(7k^3+72k^2+221k+156>0\). The remaining products are nonnegative. \end{proof}

\begin{lemma}[bulk widths four, five, and six]\label{res:A.32} Suppose \(t\in\{4,5,6\}\). Then \eqref{eq:capacity-7} holds except possibly at
\[
(r,s,t)=(4,4,4),(5,4,4),(6,4,4),(7,4,4),
(5,5,4),(6,5,4),(6,6,4),(5,5,5).
\]
\end{lemma}
\begin{proof} Let \(\Gamma=[|C|-c(t)]_+\). Since every plane excess is at least one, \eqref{eq:capacity-6} and Lemma~\ref{res:A.30} give
\(m\ge M+\Gamma+3\).
The exact row sum \eqref{eq:capacity-2}, monotone in \(r\), gives the following bounds:
\[
\begin{array}{c|c}
t&\text{bound on }\Gamma\\\hline
4&\Gamma\ge6s-14\ (r=s),\quad \Gamma\ge6s-12\ (r=s+1),\quad \Gamma\ge6s-11\ (r\ge s+2)\\
5&\Gamma\ge10s-28,\quad \Gamma\ge10s-24\text{ if }r\ge s+1\\
6&\Gamma\ge15s-52.
\end{array}
\]
For \(t=4\), this follows from \eqref{eq:capacity-3} and \(c(4)=1\). For \(t=5\), there are ten rows \(y,z>0,y+z\le5\). At \(r=s\) their lengths are \(s-\max(y,z)\), and \(\sum\max(y,z)=27\). At \(r=s+1\) their lengths are \(s-\max(y-1,z)\), whose subtracted sum is 23. For \(t=6\), the fifteen rows at \(r=s\) give subtracted sum \(\sum\max(y,z)=48\). These sums can be evaluated layer by layer for \(y+z=2,3,4,5,6\): the sums of \(\max(y,z)\) are respectively \(1,4,8,14,21\); the first four sum to 27 and all five to 48. For \(\max(y-1,z)\), the first four layer sums are \(1,3,7,12\), summing to 23. Since \(s\ge t\), the row lengths are nonnegative. Monotonicity in \(r\) preserves the estimates. Subtracting \(c(5)=1\) and \(c(6)=4\), respectively, gives the bounds for \(t=5,6\).

The square difference simplifies to
\[
4(M+\Gamma+3)-(r+s+t+2)^2
=(\delta-t+3)^2-4st+4\Gamma+3.
\]
For \(t=4,\delta=0\) the lower bound is \(8s-52\), nonnegative for \(s\ge7\). For \(\delta=1\) it is \(8s-45\), nonnegative for \(s\ge6\). For \(\delta\ge2,s\ge5\), it is \(\delta(\delta-2)+8s-40\ge0\). For \(s=4,\delta\ge4\), it is \((\delta-4)(\delta+2)\ge0\). The excluded cases in width four are exactly those not covered by these alternatives.

For \(t=5,\delta=0,s\ge6\), the bound is \(20s-105\ge0\). For \(\delta\ge1,s\ge5\), it is \((\delta-2)^2+20s-93\ge0\). Only \((5,5,5)\) remains. For \(t=6\) the bound is \((\delta-3)^2+36s-205\ge0\) for every \(s\ge6\). In every covered case the square is at most \(4m\), proving the ceiling inequality. \end{proof}

The remaining arguments use small explicit arithmetic tables. To make every entry reproducible, the functions used in the tables are exactly the finite sums and minima of Lemma~\ref{res:A.8}. In particular
\[
Q(r;a,b)=\sum_{x=1}^{a}\min\{b,[x+b-r]_+\},
\]
and \(F\) is computed from \(Q,W,\Phi\). The formula for \(Q\) counts the integers \(y\in[1,b]\) with \(y>r-x\). Thus the tables require only integer addition, multiplication, comparisons, and loops over the explicitly displayed finite intervals. The preceding estimates reduce the remaining arithmetic to the displayed finite domains.

\begin{lemma}[the eight remaining width profiles]\label{res:A.33}\label{res:A.34}
The full-support capacity \eqref{eq:capacity-7} holds at all eight triples left by
Lemma~\ref{res:A.32}.
\end{lemma}
\begin{proof}
For each displayed triple put
\[
 \Lambda=\left\lceil(r+s+t+2)^2/4\right\rceil,\qquad
 M=1+3r+\binom r2+\binom s2+\binom t2,\qquad \Delta=\Lambda-M.
\]
Here $t=4$ or $5$, so $c(t)=1$ in Lemma~\ref{res:A.30}.  Put $\ell_{\mathrm{core}}=|C|-1$;
then $|H|\ge \ell_{\mathrm{core}}$.
If $A\ge7$, rectangle excess on the two incident planes (one point
on each diagonal above its threshold), and a unit on the third, give
\[
 F(r;A,B)+F(s;A,C_0)+F(t;B,C_0)
 \ge(A+B-r)+(A+C_0-s)+1\ge17-r-s.
\]
For $B\ge7$ or $C_0\ge7$, the corresponding bounds are
$17-r-t$ or $17-s-t$, both at least $17-r-s$.
The table below has $\Delta+r+s\le \ell_{\mathrm{core}}+17$ in every row.
Consequently $m\ge M+\ell_{\mathrm{core}}+17-r-s\ge \Lambda$ whenever any coordinate of $g$
is at least seven.

It remains to check the explicitly finite domain
\[
 \Omega_r=\{(a,b,c)\in\{1,\ldots,6\}^3:a+b+c\ge r+2\},
\]
which contains every remaining $g$ by terminal capacity.
The exact interior resource at $g=(a,b,c)$ is
\[
 h_{r,s,t}(a,b,c)
 =abc-1+\#\{q\in C:(a,b,c)\nleq q,\ q\nleq(a,b,c)\}.
\]
Indeed, the positive proper box contributes $abc-1$ points, and its
union with the punctured core adds precisely the core points
incomparable with $g$.  Let $\mathcal B_{r,s,t}(a,b,c)$ be the full
right side $\mathsf{BF}$ of \eqref{eq:capacity-6}, with this exact value in place of
$|H|$.  The following table gives its minimum excess over $\Lambda$.

% CAPACITY_V03_EIGHT_PROFILE_TABLE_BEGIN
\begin{center}
\begin{tabular}{c|r|r|r|r}
 $(r,s,t)$&$|C|$&$\Delta$&$|\Omega_r|$
    &$\min_{\Omega_r}(\mathcal B-\Lambda)$\\\hline
 $(4,4,4)$&11&18&206&0\\
 $(5,4,4)$&13&19&196&3\\
 $(6,4,4)$&14&18&181&6\\
 $(7,4,4)$&14&18&160&9\\
 $(5,5,4)$&17&22&196&2\\
 $(6,5,4)$&19&23&181&5\\
 $(6,6,4)$&23&26&181&4\\
 $(5,5,5)$&23&27&196&5
\end{tabular}
\end{center}
\hypertarget{tab:capacity-eight-profiles}{}
% CAPACITY_V03_EIGHT_PROFILE_TABLE_END

Every table entry is a finite calculation from the definitions:
enumerate $1\le a,b,c\le6$ with $a+b+c\ge r+2$, count the displayed
subset of the fixed core in $h$, evaluate $F,L_i,J$ by Lemmas~\ref{res:A.8}--\ref{res:A.9},
and retain the least $\mathcal B-\Lambda$.
The eight domains contain $1497$ triples in total.
All minima are nonnegative, proving the capacity on the entire finite
domain.  Together with the preceding unbounded-coordinate estimate,
this proves every displayed profile. \end{proof}

\begin{theorem}[ordered full-support capacity for widths at least four]\label{res:A.35} Every ordered cyclic staircase with a full boundary and \(t\ge4\) satisfies \eqref{eq:capacity-7}.
\end{theorem}
\begin{proof} For \(t\ge7\) use Lemma~\ref{res:A.31}. For \(4\le t\le6\), Lemma~\ref{res:A.32} covers all but its eight explicit triples; Lemma~\ref{res:A.33} covers all eight. \end{proof}


\begin{proof}[Completion of Theorem~\ref{thm:capacity}]
Without full support use Lemma~\ref{res:A.22}. With full support,
widths zero, one, two, three, and at least four are covered by
Lemma~\ref{res:A.19}, Lemma~\ref{res:A12.1},
Theorem~\ref{res:A12.8}, Theorem~\ref{res:A3.9}, and
Theorem~\ref{res:A.35}, respectively.
Thus \eqref{eq:capacity-7}--\eqref{eq:capacity-8} hold in every ordered case. Coordinate transport
gives them for arbitrary pair counts. By Lemma~\ref{res:A.2}, their
common conclusion is $(|B|-2)^2\le4m$, hence
$|B|\le2+\lfloor\sqrt{4m}\rfloor$.
\end{proof}
\section{Consequences}\label{sec:betti}
\subsection{Total Betti numbers}

For a field $\Bbbk$, put $R=\Bbbk[S]$ and
$P=\Bbbk[x_0,x_1,x_2,x_3]$, graded by $\deg x_i=g_i$, with $x_i\mapsto t^{g_i}$ and
$(g_0,g_1,g_2,g_3)=(m,n_1,n_2,n_3)$.  We write
$\beta_i=\dim_{\Bbbk}\operatorname{Tor}_i^P(R,\Bbbk)$ for the
\emph{total} Betti numbers, and $\tau=|\operatorname{PF}(S)|$ for the type, where $\operatorname{PF}(S)$ denotes the pseudo-Frobenius numbers.
Thus $\beta_i$ sums over all internal degrees; it is a total rank,
whereas a graded Betti number records a specified internal degree.
The number of distinct Betti degrees is a different invariant.

For a finite downset, $\max D$ denotes its set of maximal elements
in the coordinatewise order. We record a consequence of the standard
Scarf-deformation construction, with the planar count included to show
the dependence on full-support generators. General background on
three-dimensional staircases and generic monomial ideals is given in
\cite[Chapters~3 and~6]{MillerSturmfels05}.

\begin{lemma}[Staircase socle]\label{lem:socle}
Let $D\subset\mathbb N^3$ be a nonempty finite downset, let
$b_D=|\partial D|$, let $f$ count the points of $\partial D$ with all
coordinates positive, and let $v=|\max D|$.  Then
\[
v\le b_D+f-2.
\]\end{lemma}
\begin{proof}
For $Q=\Bbbk[x,y,z]$, the monomial ideal
$J=(x^{b_1}y^{b_2}z^{b_3}:b\in\partial D)$ is Artinian, has $b_D$ minimal
generators, and $Q/J$ has socle basis indexed by $\max D$.
Perturb the positive exponents of the generators generically, preserving
zero exponents and every strict comparison between old exponents.
The resulting ideal remains Artinian with the same supports.  Its Scarf
complex supports a free resolution of $Q/J$ after restoring the original
labels, by the deformation construction of
\cite[Theorem~4.3 and Corollary~4.4]{BPS98}.
By the radial projection of \cite[Construction~5.4 and Corollary~5.5]{BPS98}, this complex triangulates a triangle, with support-zero vertices on the corresponding boundary faces.
Its $f$ full-support vertices are interior, and its other $b_D-f$ vertices
are on the boundary.  If $\nu_1$ and $\nu_2$ denote its numbers of edges and
triangles, Euler's formula and edge incidence give
$b_D-\nu_1+\nu_2=1$ and $3\nu_2=2\nu_1-(b_D-f)$.  Thus $\nu_2=b_D+f-2$.
The last free module has rank $\nu_2$, whereas the minimal last rank is the
socle dimension $v$.  Hence $v\le \nu_2$. \end{proof}

\medskip
\begin{proposition}[Simultaneous Betti bounds]\label{prop:betti}
Let $S$ have embedding dimension four and multiplicity $m$, and let
$D$ be its canonical Ap\'ery staircase.  Put $b_D=|\partial D|$,
$v=|\max D|$, and let $\epsilon\in\{0,1\}$ record the existence of a
full-support border point.  Then
\[
(\beta_0,\beta_1,\beta_2,\beta_3)
 = (1,\eta(S),\eta(S)+\tau-1,\tau),
\qquad
\tau\le v\le b_D+\epsilon-2.
\]
Consequently, with $T(m)=2+\lfloor\sqrt{4m}\rfloor$,
\[
\beta_1\le T(m),\qquad
\beta_3\le T(m)-1,\qquad
\beta_2\le 2T(m)-2.
\]
If $\epsilon=0$, the last two bounds improve to
$\beta_3\le T(m)-2$ and $\beta_2\le2T(m)-3$.\end{proposition}

\begin{proof}
The one-dimensional domain $R$ is Cohen--Macaulay, so its minimal
resolution over $P$ has length three.  Its first rank is $\eta(S)$,
its last rank is $\tau$, and its alternating rank sum is zero; these
standard identities give the displayed sequence
\cite[Section~1]{MSSurvey25}.
The quotient $A=R/(t^m)$ satisfies $\dim_{\Bbbk}\operatorname{Soc}(A)=\tau$, since $t^m$ is regular on $R$. It has basis
$\{t^{w(q)}:q\in D\}$.
Its socle is spanned by those basis elements for which
$w(q)+n_1,w(q)+n_2,w(q)+n_3$ all lie outside $\operatorname{Ap}(S,m)$.
Each such $q$ is coordinatewise maximal in $D$, proving $\tau\le v$.
The preceding lemma and the full-support uniqueness of Lemma~\ref{res:A.1} give
$v\le b_D+\epsilon-2$.  Finally the capacity theorem gives $b_D\le T(m)$,
while the border presentation gives $\eta(S)\le b_D$.
Substitution proves the claims. \end{proof}

\begin{remark}[Betti sequences realized by the known families]\label{rem:betti-family}
For every $q\ge6$ and $4m\ge(q-2)^2$, the semigroups constructed in
\cite[Theorem~6.3]{EHOOPK24} realize
\[
(\beta_0,\beta_1,\beta_2,\beta_3)=(1,q,2q-4,q-3).
\]
Indeed, the published Ap\'ery descriptions have unique factorizations.
For $q=2k+4$, their exponent region is the union of two coordinate
triangles of side $k$ and a possible extension of one axis; it has
$2k+1$ maximal points.  For $q=2k+3$, the two triangle sides are $k$
and $k-1$, with the same kind of axis extension; there are $2k$
maximal points. An extended axis replaces its former maximal point, so these counts also hold when the extension is empty. Uniqueness of Ap\'ery factorizations makes these
points precisely the socle basis of $R/(t^m)$, so the type is $q-3$.
The total-rank identity gives the other entries.
\end{remark}

\begin{remark}[Extremal windows]\label{rem:extremal-windows}
For fixed \(m\ge4\), let \(\mathcal S_4(m)\) be the class of numerical
semigroups of multiplicity \(m\) and embedding dimension four, and put
\[
\mathcal T_4(m)=\max_{S\in\mathcal S_4(m)}\tau(S),
\qquad
\mathcal B_{2,4}(m)=\max_{S\in\mathcal S_4(m)}\beta_2(S).
\]
The class is nonempty by Theorem~\ref{thm:main}, and
Proposition~\ref{prop:betti} bounds both integer-valued invariants, so
these maxima exist. The published families at \(q=T(m)\) give
\[
T(m)-3\le\mathcal T_4(m)\le T(m)-1,
\qquad
2T(m)-4\le\mathcal B_{2,4}(m)\le2T(m)-2.
\]
\end{remark}

\subsection{The attainable spectrum}
\medskip
\begin{corollary}[Complete attainable spectrum]\label{cor:spectrum}
Let \(\mathcal A_4(m)\) denote the set of minimal-presentation cardinalities
of numerical semigroups of multiplicity \(m\) and embedding dimension four.
Then
\[
\mathcal A_4(m)=
\begin{cases}
\{6\},&m=4,\\
\{5,6\},&m=5,\\
\{4,5,6\},&m=6,\\
\{4,5,6,7\},&m=7,\\
\{3,4,\ldots,T(m)\},&m\ge8.
\end{cases}
\]
In particular, Conjecture~7.1 of \cite{EHOOPK24} holds.\end{corollary}

\begin{proof}
The standard inequality \(\eta(S)\ge e(S)-1\), with equality precisely
for complete intersections, gives \(\eta(S)\ge3\). By
\cite[Propositions~6.2 and~7.6]{EHOOPK24}, cardinality three occurs
precisely when \(m\ge8\). Their Theorem~3.4 gives
\(\eta(S)\ge\binom42-(m-4)=10-m\), so cardinalities four and five
require \(m\ge6\) and \(m\ge5\), respectively. Their Theorem~4.2,
with \(e=4\) and \(s=2\) or \(s=1\), supplies these values throughout
the respective ranges.

For \(q\ge6\), the universal upper bound gives the necessary condition
\(q\le T(m)\), equivalently \((q-2)^2\le4m\); their Theorem~6.3
gives sufficiency. These cases exhaust all cardinalities and yield the
displayed sets.
\end{proof}

\subsection{Comparison with the Macaulay bound}\label{sec:macaulay}

We use the standard Macaulay raising operation \(s\mapsto s^{\langle r\rangle}\),
with \(0^{\langle r\rangle}=0\). If
\(s=\sum_i\binom{a_i}{i}\) is its canonical degree-\(r\) expansion, then
\[
s^{\langle r\rangle}-s=\sum_i\binom{a_i}{i+1}\ge0.
\]
Specializing \cite[equation~(1)]{MS25}, choose \(r\ge2\) so that
\[
\binom{r+2}{3}\le m<\binom{r+3}{3},
\qquad s=m-\binom{r+2}{3}.
\]
Then
\[
\mathcal C(3,m)=\binom{r+2}{2}+s^{\langle r\rangle}-s
\ge\binom{r+2}{2}.
\]

\medskip
\begin{corollary}[Macaulay equality]\label{cor:macaulay}
For \(m\ge4\), equality \(\mathcal R(3,m)=\mathcal C(3,m)\) holds
if and only if \(m\le9\).\end{corollary}

\begin{proof}
Equality for \(4\le m\le9\) is \cite[Theorem~3.3]{MS25} with
\(e=3\). For \(10\le m\le15\), we have \(r=3\), so
\(\mathcal C(3,m)\ge10>T(m)\). For \(16\le m\le19\), we have
\(r=3\) and \(s\ge6\). The first top in the degree-three expansion
of \(s\) is at least four, whence
\(s^{\langle3\rangle}-s\ge\binom44=1\). Thus
\(\mathcal C(3,m)\ge11>T(m)=10\).

Finally, suppose \(m\ge20\), so \(r\ge4\), and write
\(K_r=\binom{r+2}{2}\), \(U=\binom{r+3}{3}=K_r(r+3)/3\).
Since
\[
3K_r-4r-24=\frac{3r^2+r-42}{2}>0 \qquad(r\ge4),
\]
we obtain
\[
(K_r-2)^2=K_r(K_r-4)+4>
\frac{4K_r(r+3)}3=4U>4m.
\]
As \(K_r-2>0\), this gives \(T(m)<K_r\le\mathcal C(3,m)\).
\end{proof}

\section*{Formalization and reproducibility}
The companion Lean development, with pinned dependencies and execution
records, is archived in software version~0.2.0
(\url{https://doi.org/10.5281/zenodo.24346459}). The companion modules \texttt{PaperPresentation.lean} and
\texttt{PaperTheorems.lean} provide the declarations in the first
three rows of the table. The project uses Lean~4.34.0-rc1 with Mathlib at
\begin{center}\small
\texttt{de5ce8a9a66a4aa68a9bdbb35b63a06d34d9ca11}.
\end{center}

The first three declarations are in the namespace \texttt{NSP3.Paper}.

\begin{center}\small
\begin{tabular}{@{}>{\raggedright\arraybackslash}p{.41\linewidth}>{\raggedright\arraybackslash}p{.53\linewidth}@{}}\toprule
Mathematical assertion&Lean declaration\\\midrule
Extremal presentation cardinality&\texttt{max\_presentation\_cardinality}\\
Complete attainable spectrum&\texttt{attainable\_presentation\_cardinalities}\\
Macaulay equality range&\texttt{macaulay\_equality\_iff}\\
Total Betti identities and higher vanishing&\texttt{NSP3.Betti.identities}\\
Tor finiteness and dimension identification&\texttt{NSP3.Betti.finite}, \texttt{dimension\_eq\_total}\\
Simultaneous and no-full-support bounds&\texttt{NSP3.Betti.upper}, \texttt{noFull\_upper}\\
Known-family tuples and field independence&\texttt{NSP3.Betti.realized}, \texttt{field\_independent}\\
Attained extremal windows&\texttt{NSP3.Betti.maximumSecond\_window}, \texttt{maximumThird\_window}\\\bottomrule
\end{tabular}
\end{center}
In the last five rows, unqualified names use the same
\texttt{NSP3.Betti} namespace. The Betti invariants are defined from
categorical Tor for the polynomial presentation of the semigroup
algebra. The coefficient field can have any characteristic.

Formal proofs may follow different arguments: a direct combinatorial
argument replaces the Scarf-deformation proof of Lemma~\ref{lem:socle}.
The Python scripts reproduce the displayed finite arithmetic; their
executions are not assumptions of the Lean proofs. The archive records
the proved interfaces and their proof routes.

The supplementary build records cover the 269 project modules needed
for the paper checks. They were compiled from source with pinned dependency
artifacts. The appendix module was checked against those outputs.
The reported axiom dependencies are Lean's standard axioms.

\section*{Appendix. Lean definitions and statement of Theorem~\ref{thm:main}}
In the following excerpts from \texttt{NSP3.Paper},
\texttt{AtomCoefficients S} consists of finitely supported vectors of
nonnegative coefficients on the minimal generators of $S$.
Their weight is the sum of the generator values with those coefficients.
The predicate \texttt{IsAtomPresentation S R} says that the additive
congruence generated by the finite relation set $R$ is exactly the
kernel congruence of this weight map.

\begin{quote}\small\ttfamily
\begin{tabular}{@{}l@{}}
\hspace*{0em}noncomputable def presentationCardinality (S : Set \(\mathbb N\)) :\kern.32em\(\mathbb N\) :=\\
\hspace*{1em}sInf \{n :\kern.32em\(\mathbb N\) | \(\exists\) R : Finset (AtomCoefficients S \(\times\) AtomCoefficients S),\\
\hspace*{2em}IsAtomPresentation S R \(\land\) R.card = n\}\\
\strut\\
\hspace*{0em}def presentationCardinalities (m :\kern.32em\(\mathbb N\)) : Set \(\mathbb N\) :=\\
\hspace*{1em}\{n | \(\exists\) S : AddSubmonoid \(\mathbb N\), ((S : Set \(\mathbb N\))\({}^{\mathrm c}\)).Finite \(\land\)\\
\hspace*{2em}semigroupEmbeddingDimension (S : Set \(\mathbb N\)) = 4 \(\land\)\\
\hspace*{2em}IsSemigroupMultiplicity (S : Set \(\mathbb N\)) m \(\land\)\\
\hspace*{2em}presentationCardinality (S : Set \(\mathbb N\)) = n\}
\end{tabular}
\end{quote}

Proposition~\ref{prop:border} gives a finite presentation, so the set
of cardinalities in \texttt{sInf} is nonempty. Its least element is
$\eta(S)$, since Lean proves that every inclusion-minimal presentation
of a four-generated numerical semigroup has this cardinality.
The embedding dimension counts the minimal generators, and the
multiplicity predicate says that $m$ is the least positive element.
The second definition collects the values of $\eta(S)$ at embedding
dimension four and multiplicity $m$.

\begin{quote}\small\ttfamily
\begin{tabular}{@{}l@{}}
\hspace*{0em}theorem max\_presentation\_cardinality (m :\kern.32em\(\mathbb N\)) (hm : 4 \(\le\) m) :\\
\hspace*{2em}IsGreatest (presentationCardinalities m)\\
\hspace*{3em}\(\lfloor\)(2 :\kern.32em\(\mathbb R\)) + 2 * Real.sqrt (m :\kern.32em\(\mathbb R\))\(\rfloor_+\)
\end{tabular}
\end{quote}

The predicate \texttt{IsGreatest} asserts both membership of the displayed
value in this set and that it bounds the set above. The right side equals
$2+\lfloor\sqrt{4m}\rfloor$.

\section*{Declarations}
AI tools (ChatGPT and Codex, OpenAI) contributed substantially to
mathematical research, proof construction, Lean formalization,
verification and writing. Claude (Anthropic, Opus~5.5 with high effort)
reviewed the manuscript for readability, exposition, attribution, and
the presentation of the Lean formalization.
AI-agent checks were author-side preparation,
not independent human peer review. The author takes responsibility for the mathematical claims, attribution
and final text.

The author received no funding and declares no competing interests.
No empirical data were collected. Original code in the supplementary
materials is licensed under Apache License~2.0; dependencies retain
their respective licences.

\begin{thebibliography}{99}

\bibitem{MS25} Moscariello, A., Sammartano, A.: On minimal presentations of numerical monoids. \emph{Bull. Lond. Math. Soc.} \textbf{57}(3), 878--894 (2025). \url{https://doi.org/10.1112/blms.70005}.

\bibitem{MSSurvey25} Moscariello, A., Sammartano, A.: Open problems on relations of numerical semigroups. In: \emph{Recent Progress in Ring and Factorization Theory}, Springer Proceedings in Mathematics \& Statistics, vol. 477, pp. 365--380 (2025). \url{https://doi.org/10.1007/978-3-031-75326-8_16}.

\bibitem{EHOOPK24} Elmac\i{}o\u{g}lu, C., Hilmer, K., O'Neill, C., Okandan, M., Park-Kaufmann, H.: On the cardinality of minimal presentations of numerical semigroups. \emph{Algebr. Comb.} \textbf{7}(3), 753--771 (2024). \url{https://doi.org/10.5802/alco.354}.

\bibitem{Chomicz26A} Chomicz, K.: On numerical semigroups with embedding dimension four. arXiv preprint arXiv:2604.25653v3 (2026). \url{https://arxiv.org/abs/2604.25653v3}.

\bibitem{Chomicz26B} Chomicz, K.: The type and cardinality of minimal presentations of numerical semigroups with embedding dimension four. arXiv preprint arXiv:2609.04000v1 (2026). \url{https://arxiv.org/abs/2609.04000v1}.

\bibitem{BMO25} Brower, C., McDonough, J., O'Neill, C.: Numerical semigroups, polyhedra, and posets IV: walking the faces of the Kunz cone. \emph{Algebr. Comb.} \textbf{8}(3), 597--617 (2025). \url{https://doi.org/10.5802/alco.425}.
\bibitem{BPS98}
Bayer, D., Peeva, I., Sturmfels, B.:
Monomial resolutions. \emph{Math. Res. Lett.} \textbf{5},
31--46 (1998).
\url{https://doi.org/10.4310/MRL.1998.v5.n1.a3}.

\bibitem{Bresinsky75}
Bresinsky, H.: On prime ideals with generic zero $x_i=t^{n_i}$.
\emph{Proc. Amer. Math. Soc.} \textbf{47}(2), 329--332 (1975).
\url{https://doi.org/10.1090/S0002-9939-1975-0389912-0}.

\bibitem{MillerSturmfels05}
Miller, E., Sturmfels, B.: \emph{Combinatorial Commutative Algebra}.
Graduate Texts in Mathematics, vol.~227. Springer (2005).
\url{https://doi.org/10.1007/b138602}.
\bibitem{Herzog70}
Herzog, J.: Generators and relations of abelian semigroups and semigroup rings.
\emph{Manuscripta Math.} \textbf{3}, 175--193 (1970).
\url{https://doi.org/10.1007/BF01273309}.

\bibitem{HerzogStamate14}
Herzog, J., Stamate, D.I.: On the defining equations of the tangent cone of
a numerical semigroup ring. \emph{J. Algebra} \textbf{418}, 8--28 (2014).
\url{https://doi.org/10.1016/j.jalgebra.2014.07.008}.

\bibitem{Apery46}
Ap\'ery, R.: Sur les branches superlin\'eaires des courbes alg\'ebriques.
\emph{C. R. Acad. Sci. Paris} \textbf{222}, 1198--1200 (1946).

\bibitem{RG09}
Rosales, J.C., Garc\'ia-S\'anchez, P.A.: \emph{Numerical Semigroups}.
Developments in Mathematics, vol.~20. Springer (2009).
\url{https://doi.org/10.1007/978-1-4419-0160-6}.
\end{thebibliography}
\end{document}
